% =====================================================================
\chapter{Soluciones de los ejercicios}
\label{ap:soluciones}
% =====================================================================

Cada solución se presenta en cinco partes: datos, planteamiento, desarrollo, resultado y comprobación. En los ejercicios de demostración ($\star\star\star$), los pasos del desarrollo, leídos en orden, forman la demostración redactada tal como debería entregarse. La comprobación indica una forma independiente de verificar el resultado (con frecuencia, derivar la primitiva obtenida o comparar con un valor numérico) o señala dónde se ha usado cada hipótesis. Los valores decimales se redondean a cuatro cifras significativas salvo que se indique otra cosa.

% ---------------------------------------------------------------------
\section{Tema 1}

\begin{solucion}{eje:1-inmediatas}
\datos seis integrandos elementales.
\planteamiento linealidad de la integral indefinida y \cref{tab:inmediatas}.
\pasos
\begin{enumerate}[label=(\alph*)]
  \item $\int(2x^{4}+3x^{2}+x+8)\dif x=\frac25x^{5}+x^{3}+\frac12x^{2}+8x+C$.
  \item $\sqrt x+x^{-1/2}$ tiene primitiva $\frac23x^{3/2}+2x^{1/2}+C$ en $(0,\infty)$.
  \item $\int e^{7x}\dif x=\frac17e^{7x}+C$.
  \item Con la última fila de la tabla y $f(u)=1/u$: $\int\frac{\dif x}{3x+2}=\frac13\ln\abs{3x+2}+C$ en cada intervalo que no contenga $x=-\frac23$.
  \item $\int\sin(5x)\dif x=-\frac15\cos(5x)+C$.
  \item $\int\frac{\dif x}{1+x^{2}}=\arctan x+C$.
\end{enumerate}
\resultado las primitivas anteriores.
\comprobacion derivando: (a) $2x^{4}+3x^{2}+x+8$; (b) $x^{1/2}+x^{-1/2}$; (c) $e^{7x}$; (d) $\frac{3}{3(3x+2)}$; (e) $\sin5x$; (f) $\frac{1}{1+x^{2}}$.
\end{solucion}

\begin{solucion}{eje:1-vf}
\planteamiento contrastar cada afirmación con la \cref{prop:primitivas-constante}, la \cref{def:sumas} y la \cref{obs:1-errores}.
\pasos
\begin{enumerate}[label=(\alph*)]
  \item \textbf{Falsa.} $\R\setminus\{0\}$ no es un intervalo. $F(x)=\ln\abs x$ y $G(x)=\ln\abs x+H(x)$, con $H=0$ en $(-\infty,0)$ y $H=1$ en $(0,\infty)$, son primitivas de $1/x$ y $F-G$ no es constante.
  \item \textbf{Falsa.} $\int_0^{2\pi}\sin x\dif x=\eval{-\cos x}{0}{2\pi}=0$ (\cref{tema:2}), mientras que el área vale~4 (\cref{ej:area-seno}): la integral es un área con signo.
  \item \textbf{Verdadera.} En cada subintervalo, $m_i\le M_i$ y $\Delta x_i>0$.
  \item \textbf{Falsa.} $f(x)=x$ en $[-1,1]$ es integrable, no nula y $\int_{-1}^{1}x\dif x=0$ (\cref{prop:par-impar}).
\end{enumerate}
\resultado (a) F; (b) F; (c) V; (d) F.
\comprobacion (d) es cierta si además $f\ge0$ y $f$ es continua (\cref{eje:1-nula}).
\end{solucion}

\begin{solucion}{eje:1-sumas-x3}
\datos $f(x)=x^{3}$ en $[0,1]$; $n=4$; $\Delta x=\frac14$.
\planteamiento evaluar $f$ en los extremos izquierdos $0,\frac14,\frac12,\frac34$, en los derechos $\frac14,\dots,1$ y en los puntos medios $\frac18,\frac38,\frac58,\frac78$.
\pasos
\begin{enumerate}[label=\arabic*.]
  \item Izquierda: $\frac14\bigl(0+\frac{1}{64}+\frac{8}{64}+\frac{27}{64}\bigr)=\frac14\cdot\frac{36}{64}=\frac{9}{64}=0{,}1406$.
  \item Derecha: $\frac14\bigl(\frac{1}{64}+\frac{8}{64}+\frac{27}{64}+\frac{64}{64}\bigr)=\frac14\cdot\frac{100}{64}=\frac{25}{64}=0{,}3906$.
  \item Punto medio: $\frac14\cdot\frac{1+27+125+343}{512}=\frac14\cdot\frac{496}{512}=\frac{31}{128}=0{,}2422$.
  \item Como $f$ es creciente en $[0,1]$, el mínimo en cada subintervalo está en el extremo izquierdo y el máximo en el derecho: la suma por la izquierda es $\Lsum$ y la suma por la derecha es $\Usum$.
\end{enumerate}
\resultado $\Lsum=\frac{9}{64}$, $\Usum=\frac{25}{64}$, punto medio $\frac{31}{128}$.
\comprobacion los errores respecto a $\frac14$ son $0{,}1094$, $0{,}1406$ y $0{,}0078$: la regla del punto medio es mucho más precisa. Además $\Usum-\Lsum=\frac14\bigl(f(1)-f(0)\bigr)=\frac{16}{64}$, como predice el \cref{teo:monotona}.
\end{solucion}

\begin{solucion}{eje:1-sumas-strang}
\datos $f(x)=x^{3}$ en $[0,2]$; $\Delta x=\frac12$; puntos $0,\frac12,1,\frac32,2$.
\planteamiento $f$ es creciente: $\Lsum$ usa los extremos izquierdos y $\Usum$ los derechos.
\pasos
\begin{enumerate}[label=\arabic*.]
  \item Valores: $f=0,\ \frac18,\ 1,\ \frac{27}{8},\ 8$.
  \item $\Lsum=\frac12\bigl(0+\frac18+1+\frac{27}{8}\bigr)=\frac12\cdot\frac{36}{8}=\frac94=2{,}25$.
  \item $\Usum=\frac12\bigl(\frac18+1+\frac{27}{8}+8\bigr)=\frac12\cdot\frac{100}{8}=\frac{25}{4}=6{,}25$.
  \item Media: $\frac{2{,}25+6{,}25}{2}=4{,}25$.
\end{enumerate}
\resultado $\Lsum=2{,}25\le4\le\Usum=6{,}25$; la media, $4{,}25$, tiene error $0{,}25$, frente a $1{,}75$ y $2{,}25$.
\comprobacion la media de las dos sumas es la regla del trapecio; $\Usum-\Lsum=\frac12(8-0)=4$, de acuerdo con el \cref{teo:monotona}.
\end{solucion}

\begin{solucion}{eje:1-limite-x3}
\datos $f(x)=x^{3}$ en $[0,b]$; $x_k=\frac{kb}{n}$; $\sum_{k=1}^{n}k^{3}=\frac{n^{2}(n+1)^{2}}{4}$.
\planteamiento la suma por la derecha es $\frac bn\sum_{k=1}^{n}f(x_k)$; $f$ es continua, luego integrable, y la suma tiende a la integral (\cref{prop:sumas-riemann} y \cref{teo:continua}).
\pasos
\begin{enumerate}[label=\arabic*.]
  \item $\Rsum_n=\frac bn\sum_{k=1}^{n}\frac{k^{3}b^{3}}{n^{3}}=\frac{b^{4}}{n^{4}}\cdot\frac{n^{2}(n+1)^{2}}{4}=\frac{b^{4}}{4}\Bigl(1+\frac1n\Bigr)^{2}$.
  \item $\lim_{n\to\infty}\Rsum_n=\frac{b^{4}}{4}$.
\end{enumerate}
\resultado $\int_0^bx^{3}\dif x=\frac{b^{4}}{4}$.
\comprobacion con $b=2$ se obtiene $4$, el valor del \cref{eje:1-sumas-strang}; con $b=1$, $\frac14$.
\end{solucion}

\begin{solucion}{eje:1-limite-seno}
\datos $\Rsum_n=\frac1n\sum_{k=1}^{n}\sin\frac{kb}{n}$, $b>0$.
\planteamiento reconocer una suma de Riemann por la derecha de $g(x)=\sin(bx)$ en $[0,1]$ con la partición uniforme.
\pasos
\begin{enumerate}[label=\arabic*.]
  \item Con $x_k=\frac kn$ y $\Delta x=\frac1n$: $g(x_k)=\sin\frac{kb}{n}$, luego $\Rsum_n=\sum_{k=1}^{n}g(x_k)\Delta x$.
  \item $g$ es continua, luego integrable, y para la partición uniforme $\Usum-\Lsum\to0$ (\cref{teo:continua}); por la \cref{prop:sumas-riemann}, $\Rsum_n\to\int_0^{1}\sin(bx)\dif x$.
  \item Con la primitiva $-\frac{\cos(bx)}{b}$ y la regla de Barrow: $\int_0^{1}\sin(bx)\dif x=\frac{1-\cos b}{b}$.
\end{enumerate}
\resultado $\displaystyle\lim_{n\to\infty}\Rsum_n=\frac{1-\cos b}{b}$.
\comprobacion con $b=\pi$ el límite es $\frac2\pi\approx0{,}6366$, el valor medio de $\sin$ en $[0,\pi]$ (\cref{ej:media-seno}), como debe ser: $\Rsum_n$ es la media de $n$ valores equiespaciados del seno en $[0,\pi]$.
\end{solucion}

\begin{solucion}{eje:1-limites-sumas}
\datos dos sumas con $n$ sumandos.
\planteamiento sacar factor $\frac1n$ y escribir cada sumando como $f(\frac kn)$.
\pasos
\begin{enumerate}[label=(\alph*)]
  \item $\frac{n}{n^{2}+k^{2}}=\frac1n\cdot\frac{1}{1+(k/n)^{2}}$: suma de Riemann por la derecha de $\frac{1}{1+x^{2}}$ en $[0,1]$, que tiende a $\int_0^{1}\frac{\dif x}{1+x^{2}}=\frac\pi4$.
  \item $\frac{1}{n+k}=\frac1n\cdot\frac{1}{1+k/n}$: suma de Riemann de $\frac{1}{1+x}$ en $[0,1]$, que tiende a $\int_0^{1}\frac{\dif x}{1+x}=\ln2$.
\end{enumerate}
\resultado (a) $\frac\pi4\approx0{,}7854$; (b) $\ln2\approx0{,}6931$.
\comprobacion con $n=100$ las sumas valen $0{,}7829$ y $0{,}6907$; con $n=1000$, $0{,}7851$ y $0{,}6929$. Ambas funciones son decrecientes, y las sumas por la derecha son sumas inferiores: se acercan por debajo.
\end{solucion}

\begin{solucion}{eje:1-condicion}
\datos $f(x)=3x^{2}-\frac2x+e^{x}$ en $(0,\infty)$; $F(1)=2$.
\planteamiento hallar una primitiva y fijar la constante con la condición (\cref{prop:primitivas-constante}).
\pasos
\begin{enumerate}[label=\arabic*.]
  \item En $(0,\infty)$: $F(x)=x^{3}-2\ln x+e^{x}+C$.
  \item $F(1)=1-0+e+C=2$, luego $C=1-e$.
  \item En $\R\setminus\{0\}$, las primitivas son $x^{3}-2\ln\abs x+e^{x}+C_1$ para $x<0$ y $x^{3}-2\ln\abs x+e^{x}+C_2$ para $x>0$, con $C_1$ y $C_2$ independientes. Una condición en $x=1$ fija $C_2$ y deja libre $C_1$.
\end{enumerate}
\resultado $F(x)=x^{3}-2\ln x+e^{x}+1-e$.
\comprobacion $F'(x)=3x^{2}-\frac2x+e^{x}$ y $F(1)=1+e+1-e=2$.
\end{solucion}

\begin{solucion}{eje:1-estimar}
\datos (a) $\int_0^{1}\frac{\dif x}{1+x^{3}}$; (b) $\int_{10}^{20}\sqrt{x^{2}+1}\dif x$.
\planteamiento acotar el integrando por una función de integral conocida y aplicar la monotonía del \cref{teo:propiedades}(iii).
\pasos
\begin{enumerate}[label=(\alph*)]
  \item Para $x\in[0,1]$, $(1-x^{3})(1+x^{3})=1-x^{6}\le1$, luego $\frac{1}{1+x^{3}}\ge1-x^{3}$. Por la monotonía y el \cref{eje:1-limite-x3} con $b=1$: $\int_0^{1}\frac{\dif x}{1+x^{3}}\ge\int_0^{1}(1-x^{3})\dif x=1-\frac14=0{,}75>0{,}65$.
  \item Para todo $x$, $\sqrt{x^{2}+1}>\sqrt{x^{2}}=\abs x$. En $[10,20]$, $\int_{10}^{20}x\dif x$ es el área de un trapecio de bases $10$ y $20$ y altura $10$, es decir, $150$. Por la monotonía, $\int_{10}^{20}\sqrt{x^{2}+1}\dif x\ge150$, y la desigualdad es estricta porque la diferencia de los integrandos es continua, positiva y de integral no nula (\cref{eje:1-nula}).
\end{enumerate}
\resultado ambas estimaciones son ciertas; se usa la monotonía de la integral.
\comprobacion numéricamente, las integrales valen $0{,}8356$ y $150{,}35$.
\end{solucion}

\begin{solucion}{eje:1-exp-sumas}
\datos $f(x)=e^{x}$ en $[0,1]$; partición uniforme $x_k=\frac kn$.
\planteamiento $f$ es creciente, luego integrable (\cref{teo:monotona}) con $\Usum-\Lsum\to0$; basta calcular $\lim\Usum(f,P_n)$.
\pasos
\begin{enumerate}[label=\arabic*.]
  \item Como $f$ es creciente, $M_k=e^{k/n}$ y $\Usum_n=\frac1n\sum_{k=1}^{n}e^{k/n}$.
  \item Es una suma geométrica de razón $q=e^{1/n}$: $\sum_{k=1}^{n}q^{k}=q\frac{q^{n}-1}{q-1}=e^{1/n}\frac{e-1}{e^{1/n}-1}$.
  \item Con $h=\frac1n$: $\Usum_n=(e-1)\,e^{h}\,\frac{h}{e^{h}-1}$.
  \item Cuando $h\to0$, $e^{h}\to1$ y $\frac{e^{h}-1}{h}\to(e^{x})'\big|_{x=0}=1$. Luego $\Usum_n\to e-1$.
  \item Por el \cref{teo:monotona}, $\Usum_n-\Lsum_n=\frac1n(e-1)\to0$, y como $\Lsum_n\le\int_0^1e^x\dif x\le\Usum_n$, la integral es el límite común.
\end{enumerate}
\resultado $\int_0^{1}e^{x}\dif x=e-1\approx1{,}718$.
\comprobacion con $n=10$, $\Usum_{10}=1{,}806$, mayor que $e-1$, como debe ser una suma superior; la regla de Barrow da $\eval{e^{x}}{0}{1}=e-1$.
\end{solucion}

\begin{solucion}{eje:1-diferencia}
\datos $f(x)=x^{2}$ en $[0,b]$; $x_k=\frac{kb}{n}$.
\planteamiento como $f$ es creciente en $[0,b]$, $M_k=f(x_k)$ y $m_k=f(x_{k-1})$; la diferencia es telescópica.
\pasos
\begin{enumerate}[label=\arabic*.]
  \item $\Usum-\Lsum=\sum_{k=1}^{n}(M_k-m_k)\frac bn=\frac bn\sum_{k=1}^{n}\bigl(x_k^{2}-x_{k-1}^{2}\bigr)=\frac bn\bigl(x_n^{2}-x_0^{2}\bigr)=\frac bn\cdot b^{2}=\frac{b^{3}}{n}$.
  \item Interpretación: la diferencia es la suma de $n$ rectángulos de anchura $\frac bn$ cuyas alturas $f(x_k)-f(x_{k-1})$ suman $f(b)-f(0)=b^{2}$; desplazados horizontalmente hasta la misma columna, se apilan en un único rectángulo de base $\frac bn$ y altura $b^{2}$.
  \item Dado $\varepsilon>0$, si $n>\frac{b^{3}}{\varepsilon}$ se tiene $\Usum(f,P_n)-\Lsum(f,P_n)<\varepsilon$, y el \cref{teo:criterio} prueba que $f$ es integrable.
\end{enumerate}
\resultado $\Usum-\Lsum=\frac{b^{3}}{n}\to0$; $f$ es integrable en $[0,b]$.
\comprobacion coincide con el \cref{ej:x2}; con $b=1$ y $n=4$, $\frac{30}{64}-\frac{14}{64}=\frac{16}{64}=\frac14=\frac{b^{3}}{n}$.
\end{solucion}

\begin{solucion}{eje:1-positiva}
\datos $f$ integrable, $f\ge0$; para la segunda parte, $f\le g$ integrables.
\planteamiento acotar inferiormente las sumas inferiores y usar la linealidad.
\pasos
\begin{enumerate}[label=\arabic*.]
  \item Sea $P$ una partición. Como $f(x)\ge0$ para todo $x$, $0$ es cota inferior de $f$ en cada subintervalo, y el ínfimo, que es la mayor de las cotas inferiores, cumple $m_i\ge0$.
  \item Por tanto $\Lsum(f,P)=\sum m_i\Delta x_i\ge0$ para toda $P$, y $\int_a^bf=\sup_P\Lsum(f,P)\ge\Lsum(f,P)\ge0$.
  \item Si $f\le g$, la función $g-f$ es integrable por la linealidad (\cref{teo:propiedades}(i)) y es no negativa; por el paso~2, $\int_a^b(g-f)\ge0$.
  \item Por la linealidad, $\int_a^bg-\int_a^bf=\int_a^b(g-f)\ge0$.
\end{enumerate}
\resultado $\int_a^bf\ge0$ y $\int_a^bf\le\int_a^bg$.
\comprobacion la hipótesis $f\ge0$ se usa en el paso~1; la integrabilidad de $f$, para identificar la integral con el supremo de las sumas inferiores.
\end{solucion}

\begin{solucion}{eje:1-un-punto}
\datos $f(c)=1$ y $f(x)=0$ si $x\neq c$, en $[a,b]$.
\planteamiento probar que $\Lsum(f,P)=0$ siempre y que $\Usum(f,P)$ puede hacerse tan pequeña como se quiera; aplicar el \cref{teo:criterio}.
\pasos
\begin{enumerate}[label=\arabic*.]
  \item Todo subintervalo $[x_{i-1},x_i]$ tiene longitud positiva, luego contiene puntos distintos de $c$, en los que $f=0$. Como $f\ge0$, $m_i=0$ y $\Lsum(f,P)=0$ para toda $P$.
  \item $M_i=1$ si $c\in[x_{i-1},x_i]$ y $M_i=0$ en otro caso. El punto $c$ pertenece como mucho a dos subintervalos (dos si es un punto de la partición interior), luego $\Usum(f,P)\le2\norm{P}$.
  \item Dado $\varepsilon>0$, sea $P$ la partición uniforme con $\frac{b-a}{n}<\frac\varepsilon2$. Entonces $\Usum(f,P)-\Lsum(f,P)\le2\norm P<\varepsilon$, y $f$ es integrable.
  \item La integral está entre $\Lsum(f,P)=0$ y $\Usum(f,P)<\varepsilon$ para todo $\varepsilon>0$, luego $\int_a^bf=0$.
  \item Si $g$ es integrable y $\tilde g$ coincide con $g$ salvo en $c$, entonces $\tilde g=g+\bigl(\tilde g(c)-g(c)\bigr)f$, y por la linealidad $\int\tilde g=\int g+\bigl(\tilde g(c)-g(c)\bigr)\cdot0=\int g$.
\end{enumerate}
\resultado $f$ es integrable, $\int_a^bf=0$, y modificar una función integrable en un punto no altera su integral.
\comprobacion repitiendo el argumento con varios puntos, lo mismo vale para un número finito de modificaciones.
\end{solucion}

\begin{solucion}{eje:1-nula}
\datos $f$ continua, $f\ge0$ en $[a,b]$, $\int_a^bf=0$.
\planteamiento reducción al absurdo: si $f$ fuera positiva en un punto, lo sería en un intervalo, y la integral sería positiva.
\pasos
\begin{enumerate}[label=\arabic*.]
  \item Supóngase que existe $x_0\in[a,b]$ con $f(x_0)>0$. Por la continuidad, con $\varepsilon=f(x_0)/2$ existe $\delta>0$ tal que $f(x)>f(x_0)/2$ para todo $x\in J=[x_0-\delta,x_0+\delta]\cap[a,b]$.
  \item $J$ es un intervalo de longitud $\ell>0$ (al menos $\min\{\delta,b-a\}$).
  \item Sea $h=\frac{f(x_0)}{2}$ en $J$ y $h=0$ fuera de $J$. Entonces $0\le h\le f$ en $[a,b]$, $h$ es integrable (es constante en tres trozos) y $\int_a^bh=\frac{f(x_0)}{2}\ell>0$.
  \item Por la monotonía, $\int_a^bf\ge\int_a^bh>0$, en contra de la hipótesis. Luego $f(x)=0$ para todo $x$.
  \item Sin continuidad la conclusión falla: la función del \cref{eje:1-un-punto} es no negativa, tiene integral nula y vale~1 en $c$.
\end{enumerate}
\resultado $f\equiv0$; la continuidad es imprescindible.
\comprobacion la hipótesis $f\ge0$ se usa en el paso~3 ($h\le f$ fuera de $J$); sin ella, $f(x)=x$ en $[-1,1]$ es un contraejemplo.
\end{solucion}

\begin{solucion}{eje:1-racionales-irr}
\datos $g(x)=x$ si $x\in\Q$ y $g(x)=0$ si $x\notin\Q$, en $[0,1]$.
\planteamiento calcular $m_i$ y $M_i$ usando que racionales e irracionales son densos.
\pasos
\begin{enumerate}[label=\arabic*.]
  \item $g\ge0$ y todo subintervalo contiene irracionales, donde $g=0$: $m_i=0$ y $\Lsum(g,P)=0$ para toda $P$. Luego $\underline I=0$.
  \item En $[x_{i-1},x_i]$, $g(x)\le x\le x_i$; y hay racionales $q$ tan cerca de $x_i$ como se quiera, con $g(q)=q$. Luego $M_i=x_i$ y $\Usum(g,P)=\sum x_i\Delta x_i=\Usum(\mathrm{id},P)$, la suma superior de la identidad $\mathrm{id}(x)=x$.
  \item La identidad es creciente, luego integrable, y $\inf_P\Usum(\mathrm{id},P)=\int_0^{1}x\dif x=\frac12$ (área de un triángulo, o la \cref{prop:sumas-riemann} con $\Usum(\mathrm{id},P_n)=\frac{n+1}{2n}\to\frac12$).
  \item Por tanto $\overline I=\frac12\neq0=\underline I$, y $g$ no es integrable.
\end{enumerate}
\resultado $\underline I=0$, $\overline I=\frac12$; $g$ no es integrable.
\comprobacion como en el \cref{ej:dirichlet}, la falta de integrabilidad se debe a que $g$ oscila en todo subintervalo; aquí la oscilación $M_i-m_i=x_i$ no se hace pequeña al refinar.
\end{solucion}

% ---------------------------------------------------------------------
\section{Tema 2}

\begin{solucion}{eje:2-conceptos}
\datos $F(x)=\int_0^{x}\sqrt{1+t^{3}}\dif t$, $x\ge0$.
\planteamiento el integrando es continuo en $[0,\infty)$ (allí $1+t^{3}\ge1$); se aplica el \cref{teo:tfc1}.
\pasos
\begin{enumerate}[label=(\alph*)]
  \item $F(0)=\int_0^{0}=0$.
  \item Por el \cref{teo:tfc1}(ii), $F'(x)=\sqrt{1+x^{3}}$. Derivando con la regla de la cadena, $F''(x)=\frac{3x^{2}}{2\sqrt{1+x^{3}}}$.
  \item $F'(x)\ge1>0$: $F$ es estrictamente creciente. $F''(x)\ge0$: $F$ es convexa en $[0,\infty)$.
\end{enumerate}
\resultado $F(0)=0$, $F'=\sqrt{1+x^{3}}$, $F''=\frac{3x^{2}}{2\sqrt{1+x^{3}}}$; $F$ es creciente y convexa.
\comprobacion la primera parte del teorema fundamental da $F'$ sin conocer una primitiva elemental, que de hecho no existe.
\end{solucion}

\begin{solucion}{eje:2-error-barrow}
\datos $\int_{-1}^{1}x^{-2}\dif x$ «calculada» como $-2$.
\planteamiento revisar las hipótesis del \cref{teo:barrow}.
\pasos
\begin{enumerate}[label=\arabic*.]
  \item El integrando no está acotado en $[-1,1]$ (tiende a $+\infty$ en $0$), luego no es integrable en el sentido de Riemann.
  \item $F(x)=-\frac1x$ no está definida en $0$; no es continua en $[-1,1]$ ni primitiva de $x^{-2}$ en todo el intervalo.
  \item Si la integral existiera, sería $\ge0$ por la monotonía, porque el integrando es positivo. Un resultado negativo es imposible.
\end{enumerate}
\resultado el cálculo es erróneo: fallan la integrabilidad de $f$ y la hipótesis sobre $F$. Como integral impropia (\cref{tema:4}), $\int_0^{1}x^{-2}\dif x$ diverge, y también la integral pedida.
\comprobacion la \cref{prop:p-cero} con $p=2\ge1$ confirma la divergencia.
\end{solucion}

\begin{solucion}{eje:2-barrow}
\planteamiento buscar una primitiva en la \cref{tab:inmediatas} y aplicar la regla de Barrow.
\pasos
\begin{enumerate}[label=(\alph*)]
  \item $\int_1^{4}3x^{1/2}\dif x=\eval{2x^{3/2}}{1}{4}=2(8-1)=14$ y $\int_1^{4}2x^{-2}\dif x=\eval{-\frac2x}{1}{4}=-\frac12+2=\frac32$. Resultado: $14-\frac32=\frac{25}{2}$.
  \item $\eval{\sin x}{0}{\pi/2}=1$.
  \item $\eval{\frac{e^{2x}}{2}}{0}{1}=\frac{e^{2}-1}{2}\approx3{,}195$.
  \item $\eval{\ln(1+x)}{0}{1}=\ln2\approx0{,}6931$.
  \item Se separa en $[-1,0]$ y $[0,2]$: $\int_{-1}^{0}(-x)\dif x+\int_0^{2}x\dif x=\frac12+2=\frac52$.
\end{enumerate}
\resultado (a) $\frac{25}{2}$; (b) $1$; (c) $\frac{e^{2}-1}{2}$; (d) $\ln2$; (e) $\frac52$.
\comprobacion (e) es la suma de las áreas de dos triángulos, $\frac12\cdot1\cdot1+\frac12\cdot2\cdot2=\frac52$.
\end{solucion}

\begin{solucion}{eje:2-arco-seno}
\datos $y=\sin x$ y $y=\sin(ax)$, $a>0$.
\planteamiento un arco de $\sin(ax)$ va de un cero al siguiente: de $0$ a $\pi/a$, donde $\sin(ax)\ge0$.
\pasos
\begin{enumerate}[label=\arabic*.]
  \item $\int_0^{\pi}\sin x\dif x=\eval{-\cos x}{0}{\pi}=2$.
  \item $\int_0^{\pi/a}\sin(ax)\dif x=\eval{-\frac{\cos(ax)}{a}}{0}{\pi/a}=\frac{-\cos\pi+\cos0}{a}=\frac2a$.
\end{enumerate}
\resultado el área de un arco de $\sin x$ es $2$, y la de un arco de $\sin(ax)$ es $\frac2a$. Comprimir la gráfica horizontalmente en un factor $a$ divide el área por $a$.
\comprobacion la altura máxima no cambia y la base pasa de $\pi$ a $\pi/a$: el área debe escalar como la base.
\end{solucion}

\begin{solucion}{eje:2-parabola-eje}
\datos (a) $y=x^{2}-a$ con $a>0$ y el eje $x$; (b) $y=x^{2}$ con $y=2x+3$, y $y=\sqrt x$ con $y=x^{2}$.
\planteamiento hallar los cortes, decidir qué curva está encima y aplicar la \cref{def:area}.
\pasos
\begin{enumerate}[label=\arabic*.]
  \item (a) $x^{2}-a=0$ en $x=\pm\sqrt a$, y entre ellos $x^{2}-a<0$. Área: $\int_{-\sqrt a}^{\sqrt a}(a-x^{2})\dif x=\eval{ax-\frac{x^{3}}{3}}{-\sqrt a}{\sqrt a}=2\Bigl(a\sqrt a-\frac{a\sqrt a}{3}\Bigr)=\frac43a^{3/2}$.
  \item Con $a=4$: $\frac43\cdot8=\frac{32}{3}$.
  \item (b) $x^{2}=2x+3$ en $x=-1$ y $x=3$; en $x=0$ la recta vale $3>0$, luego está encima. Área: $\int_{-1}^{3}(2x+3-x^{2})\dif x=\eval{x^{2}+3x-\frac{x^{3}}{3}}{-1}{3}=9-\bigl(-\frac53\bigr)=\frac{32}{3}$.
  \item $\sqrt x=x^{2}$ en $x=0$ y $x=1$, y $\sqrt x\ge x^{2}$ en $[0,1]$. Área: $\int_0^{1}(\sqrt x-x^{2})\dif x=\frac23-\frac13=\frac13$.
\end{enumerate}
\resultado (a) $\frac{32}{3}$ y $\frac43a^{3/2}$; (b) $\frac{32}{3}$ y $\frac13$.
\comprobacion el área entre una parábola $y=x^{2}$ y una recta que la corta en $x_1<x_2$ vale $\frac{(x_2-x_1)^{3}}{6}$; con $x_2-x_1=4$ da $\frac{64}{6}=\frac{32}{3}$, y con $2\sqrt a$, $\frac{8a^{3/2}}{6}=\frac43a^{3/2}$.
\end{solucion}

\begin{solucion}{eje:2-derivar}
\planteamiento \cref{cor:leibniz} con integrandos continuos.
\pasos
\begin{enumerate}[label=(\alph*)]
  \item $f(u)=\sqrt u\sin u$ es continua en $[0,\infty)$ y $g(x)=x^{2}$: $F'(x)=\sqrt{x^{2}}\sin(x^{2})\cdot2x=2x\abs x\sin(x^{2})$; para $x>0$, $2x^{2}\sin(x^{2})$.
  \item $f(t)=\frac{1}{\sqrt{1-t^{2}}}$ es continua en $(-1,1)$ y $\sin x\in(-1,1)$ para $x\in(-\frac\pi2,\frac\pi2)$: $F'(x)=\frac{\cos x}{\sqrt{1-\sin^{2}x}}=\frac{\cos x}{\abs{\cos x}}=1$, pues $\cos x>0$. Como $F(0)=0$, $F(x)=x$ en ese intervalo; es decir, $\arcsin(\sin x)=x$.
  \item $F'(x)=e^{-x^{4}}\cdot2x-e^{-x^{2}}\cdot1=2x\,e^{-x^{4}}-e^{-x^{2}}$.
\end{enumerate}
\resultado (a) $2x\abs x\sin(x^{2})$; (b) $F'=1$, $F(x)=x$; (c) $2xe^{-x^{4}}-e^{-x^{2}}$.
\comprobacion en (c), $F(1)=\int_1^{1}=0$ y $F'(1)=2e^{-1}-e^{-1}=e^{-1}>0$: para $x$ algo mayor que~1 el intervalo $[x,x^{2}]$ crece y la integral se hace positiva.
\end{solucion}

\begin{solucion}{eje:2-hallar-f}
\datos (a) $\int_0^{x}f=2x(\sin x+1)$; (b) $\int_0^{x/2}f=2x(\sin x+1)$.
\planteamiento derivar ambos miembros (\cref{cor:leibniz}).
\pasos
\begin{enumerate}[label=\arabic*.]
  \item (a) $f(x)=2(\sin x+1)+2x\cos x$. En $x=\frac\pi2$: $f(\frac\pi2)=2\cdot2+\pi\cdot0=4$.
  \item (b) $f(\frac x2)\cdot\frac12=2(\sin x+1)+2x\cos x$, luego $f(\frac x2)=4(\sin x+1)+4x\cos x$.
  \item Para obtener $f(\frac\pi2)$ se toma $x=\pi$: $f(\frac\pi2)=4(0+1)+4\pi(-1)=4-4\pi$.
\end{enumerate}
\resultado (a) $f(\pi/2)=4$; (b) $f(\pi/2)=4-4\pi\approx-8{,}566$.
\comprobacion en ambos casos el miembro derecho se anula en $x=0$, como el izquierdo; si no fuera así, no existiría tal $f$.
\end{solucion}

\begin{solucion}{eje:2-valor-medio}
\planteamiento \cref{def:valor-medio} y resolver $f(c)=\bar f$.
\pasos
\begin{enumerate}[label=(\alph*)]
  \item $\bar f=\frac13\int_0^{3}x^{2}\dif x=\frac13\cdot9=3$; $c^{2}=3$ da $c=\sqrt3\approx1{,}732\in[0,3]$.
  \item $\bar f=\frac{1}{\ln2}\int_0^{\ln2}e^{x}\dif x=\frac{2-1}{\ln2}=\frac{1}{\ln2}\approx1{,}443$; $e^{c}=\frac{1}{\ln2}$ da $c=-\ln(\ln2)\approx0{,}3665\in[0;\,0{,}6931]$.
\end{enumerate}
\resultado (a) $\bar f=3$, $c=\sqrt3$; (b) $\bar f=1/\ln2$, $c=-\ln\ln2$.
\comprobacion en ambos casos $\bar f$ está entre el mínimo y el máximo de $f$: $0\le3\le9$ y $1\le1{,}443\le2$.
\end{solucion}

\begin{solucion}{eje:2-estudio-F}
\datos $F(x)=\int_1^{x}\frac{\dif t}{1+t^{2}}$.
\planteamiento \cref{teo:tfc1} para $F'$; comparación de integrandos para la cota.
\pasos
\begin{enumerate}[label=(\alph*)]
  \item El integrando es continuo en $\R$, luego $F$ está definida en todo $\R$. $F(1)=0$; para $x>1$, $F(x)>0$ (integrando positivo); para $x<1$, $F(x)=-\int_x^{1}<0$.
  \item $F'(x)=\frac{1}{1+x^{2}}>0$: estrictamente creciente, sin extremos relativos. $F''(x)=-\frac{2x}{(1+x^{2})^{2}}$, positiva si $x<0$ y negativa si $x>0$: convexa en $(-\infty,0)$, cóncava en $(0,\infty)$, con inflexión en $x=0$.
  \item Para $t\ge1$, $\frac{1}{1+t^{2}}<\frac{1}{t^{2}}$, luego $F(x)<\int_1^{x}\frac{\dif t}{t^{2}}=1-\frac1x<1$ para $x>1$.
  \item Como $\arctan'=\frac{1}{1+x^{2}}$, la regla de Barrow da $F(x)=\arctan x-\frac\pi4$. Límites: $\frac\pi2-\frac\pi4=\frac\pi4$ en $+\infty$ y $-\frac\pi2-\frac\pi4=-\frac{3\pi}{4}$ en $-\infty$. Además, $F(x)+\frac\pi4=\arctan x$ es impar: la gráfica es simétrica respecto del punto $(0,-\frac\pi4)$, que es el punto de inflexión.
\end{enumerate}
\resultado $F$ creciente en $\R$, inflexión en $(0,-\frac\pi4)$, $F(x)=\arctan x-\frac\pi4$, límites $\frac\pi4$ y $-\frac{3\pi}{4}$.
\comprobacion $\frac\pi4\approx0{,}785<1$, de acuerdo con la cota de (c).
\end{solucion}

\begin{solucion}{eje:2-limite}
\datos $f(t)=\frac{1}{\sqrt{1+t^{4}}}$, continua en $\R$.
\planteamiento (a) teorema del valor medio integral; (b) definición de derivada.
\pasos
\begin{enumerate}[label=(\alph*)]
  \item Por el \cref{teo:tvm-integral}, $\int_1^{1+h}f=f(c_h)\,h$ con $c_h$ entre $1$ y $1+h$. Entonces $\frac1h\int_1^{1+h}f=f(c_h)$, y como $c_h\to1$ y $f$ es continua, el límite es $f(1)$.
  \item $\frac1h\int_1^{1+h}f=\frac{F(1+h)-F(1)}{h}\to F'(1)=f(1)$, por el \cref{teo:tfc1}.
\end{enumerate}
\resultado $\frac{1}{\sqrt2}\approx0{,}7071$.
\comprobacion geométricamente, el área de una franja estrecha de anchura $h$ bajo la curva es aproximadamente $h\,f(1)$.
\end{solucion}

\begin{solucion}{eje:2-par-impar}
\datos $f$ continua en $\R$; $F(x)=\int_0^{x}f$, derivable con $F'=f$ (\cref{teo:tfc1}).
\planteamiento probar que una función auxiliar tiene derivada nula y vale $0$ en $0$.
\pasos
\begin{enumerate}[label=\arabic*.]
  \item (a) Sea $f$ par y $G(x)=F(x)+F(-x)$. Por la regla de la cadena, $G'(x)=f(x)-f(-x)=0$ para todo $x$.
  \item Por la \cref{prop:primitivas-constante}, $G$ es constante en $\R$, y $G(0)=2F(0)=0$. Luego $F(-x)=-F(x)$: $F$ es impar.
  \item (b) Sea $f$ impar y $H(x)=F(x)-F(-x)$. Entonces $H'(x)=f(x)+f(-x)=0$, $H$ es constante y $H(0)=0$. Luego $F(-x)=F(x)$: $F$ es par.
\end{enumerate}
\resultado si $f$ es par, $F$ es impar; si $f$ es impar, $F$ es par.
\comprobacion $f(x)=x^{2}$ (par) da $F(x)=\frac{x^{3}}{3}$ (impar); $f(x)=x$ (impar) da $F(x)=\frac{x^{2}}{2}$ (par). Como consecuencia se recupera la \cref{prop:par-impar}: $\int_{-a}^{a}f=F(a)-F(-a)$.
\end{solucion}

\begin{solucion}{eje:2-logaritmo}
\datos $L(x)=\int_1^{x}\frac{\dif t}{t}$, con $L'(x)=\frac1x$ y $L(1)=0$.
\planteamiento en cada apartado se compara la derivada de los dos miembros y se fija la constante en $x=1$.
\pasos
\begin{enumerate}[label=\arabic*.]
  \item (a) Sea $h(x)=L(1/x)+L(x)$ en $(0,\infty)$. Por la regla de la cadena, $h'(x)=\frac{1}{1/x}\cdot\bigl(-\frac{1}{x^{2}}\bigr)+\frac1x=-\frac1x+\frac1x=0$. Luego $h$ es constante, igual a $h(1)=0$, y $L(1/a)=-L(a)$.
  \item (b) Sea $P(n)$: «$L(x^{n})=n\,L(x)$ para todo $x>0$». $P(1)$ es trivial. Supóngase $P(n)$ y sea $k(x)=L(x^{n+1})-L(x^{n})-L(x)$. Entonces $k'(x)=\frac{(n+1)x^{n}}{x^{n+1}}-\frac{nx^{n-1}}{x^{n}}-\frac1x=\frac{n+1-n-1}{x}=0$, y $k(1)=0$. Luego $L(x^{n+1})=L(x^{n})+L(x)=nL(x)+L(x)=(n+1)L(x)$, que es $P(n+1)$.
  \item (c) En $[1,2]$, $\frac1t\ge\frac12$, luego $L(2)\ge\frac12>0$. Por (b), $L(2^{n})=nL(2)\ge\frac n2$. Dado $M$, sea $n>2M$; como $L$ es creciente, $L(x)\ge L(2^{n})>M$ para todo $x\ge2^{n}$. Luego $L(x)\to\infty$.
\end{enumerate}
\resultado $L(1/a)=-L(a)$, $L(a^{n})=nL(a)$ y $\lim_{x\to\infty}L(x)=\infty$.
\comprobacion con $L=\ln$: $\ln\frac12=-\ln2$ y $\ln8=3\ln2$. La cota $L(2^{n})\ge\frac n2$ es la misma que la de $H_{2^{n}}$ en el \cref{teo:armonica}: ambas comparan $\frac1t$ con escalones.
\end{solucion}

\begin{solucion}{eje:2-periodica}
\datos $f$ continua con $f(x+T)=f(x)$ para todo $x$.
\planteamiento derivar $G(a)=\int_a^{a+T}f$ respecto de $a$.
\pasos
\begin{enumerate}[label=\arabic*.]
  \item Por el \cref{cor:leibniz} con $g(a)=a+T$ y $h(a)=a$: $G'(a)=f(a+T)\cdot1-f(a)\cdot1=0$.
  \item Luego $G$ es constante en $\R$ (\cref{prop:primitivas-constante}).
  \item $\sin^{2}$ es continua y periódica de periodo $2\pi$, luego $\int_a^{a+2\pi}\sin^{2}x\dif x=\int_0^{2\pi}\sin^{2}x\dif x=\pi$.
\end{enumerate}
\resultado la integral sobre un periodo no depende de $a$; en particular vale $\pi$ para $\sin^{2}$.
\comprobacion con $a=-\pi$: $\int_{-\pi}^{\pi}\sin^{2}x\dif x=2\int_0^{\pi}\sin^{2}x\dif x=2\cdot\frac\pi2=\pi$ (\cref{prop:par-impar} y \cref{eje:3-trig}).
\end{solucion}

\begin{solucion}{eje:2-simetrica-f}
\datos $f$ continua en $[a,b]$ con $\int_a^{x}f=\int_x^{b}f$ para todo $x$.
\planteamiento derivar la identidad respecto de $x$.
\pasos
\begin{enumerate}[label=\arabic*.]
  \item Para $x\in(a,b)$, $\frac{d}{dx}\int_a^{x}f=f(x)$ y $\frac{d}{dx}\int_x^{b}f=\frac{d}{dx}\Bigl(\int_a^{b}f-\int_a^{x}f\Bigr)=-f(x)$, por el \cref{teo:tfc1}.
  \item Como las dos funciones de $x$ son iguales, sus derivadas también: $f(x)=-f(x)$, es decir, $f(x)=0$ en $(a,b)$.
  \item Por la continuidad, $f(a)=\lim_{x\to a^{+}}f(x)=0$ y, del mismo modo, $f(b)=0$.
\end{enumerate}
\resultado $f\equiv0$ en $[a,b]$.
\comprobacion la hipótesis dice que todo punto parte el área con signo en dos mitades iguales; con $x=a$ y $x=b$ se obtiene $\int_a^bf=0$, y la derivación muestra que eso solo es posible si $f$ es nula.
\end{solucion}

\begin{solucion}{eje:2-tvm-ponderado}
\datos $f,w$ continuas en $[a,b]$, $w\ge0$.
\planteamiento como en el \cref{teo:tvm-integral}, encerrar $\int fw$ entre $m\int w$ y $M\int w$.
\pasos
\begin{enumerate}[label=\arabic*.]
  \item Sea $W=\int_a^bw\ge0$. Si $W=0$, el \cref{eje:1-nula} da $w\equiv0$; ambos miembros son nulos y sirve cualquier $c$.
  \item Si $W>0$: por el teorema de Weierstrass, $f$ alcanza un mínimo $m=f(x_1)$ y un máximo $M=f(x_2)$. Como $w\ge0$, $m\,w(x)\le f(x)\,w(x)\le M\,w(x)$ para todo $x$.
  \item Integrando (monotonía y linealidad): $mW\le\int_a^bfw\le MW$, es decir, $\mu=\frac1W\int_a^bfw\in[f(x_1),f(x_2)]$.
  \item Por el teorema de los valores intermedios existe $c$ entre $x_1$ y $x_2$ con $f(c)=\mu$, y entonces $\int_a^bfw=f(c)\,W$.
  \item Si $w$ cambia de signo puede fallar: en $[-1,1]$ con $f(x)=w(x)=x$, $\int fw=\int_{-1}^{1}x^{2}\dif x=\frac23$, pero $f(c)\int w=f(c)\cdot0=0$ para todo $c$.
\end{enumerate}
\resultado existe $c\in[a,b]$ con $\int_a^bfw=f(c)\int_a^bw$; la hipótesis $w\ge0$ es necesaria.
\comprobacion con $w\equiv1$ se recupera el \cref{teo:tvm-integral}; la hipótesis $w\ge0$ se usa en el paso~2, al multiplicar las desigualdades por $w(x)$.
\end{solucion}

% ---------------------------------------------------------------------
\section{Tema 3}

\begin{solucion}{eje:3-elegir}
\planteamiento aplicar la estrategia de la sección~3.5 a cada integrando.
\pasos
\begin{enumerate}[label=(\alph*)]
  \item Cambio $u=x^{2}$: aparece $x\dif x=\frac12\dif u$. Resultado $\frac12\sin(x^{2})+C$.
  \item Por partes con $u=x$, $v'=\cos x$.
  \item Fracciones simples: $x^{2}-5x+6=(x-2)(x-3)$.
  \item Potencia impar del coseno: se aparta un $\cos x$ y se cambia $u=\sin x$.
  \item Sustitución trigonométrica $x=3\sin t$.
  \item Cambio $u=e^{x}$: queda $\int\frac{\dif u}{1+u^{2}}=\arctan(e^{x})+C$.
  \item Por partes con $u=\ln x$, $v'=x^{3}$.
  \item El numerador es la derivada del denominador: cambio $u=x^{2}+x+7$, resultado $\ln(x^{2}+x+7)+C$.
\end{enumerate}
\resultado las técnicas indicadas.
\comprobacion en (a), (f) y (h) la primitiva se comprueba derivando: $x\cos(x^{2})$, $\frac{e^{x}}{1+e^{2x}}$ y $\frac{2x+1}{x^{2}+x+7}$.
\end{solucion}

\begin{solucion}{eje:3-sustitucion}
\planteamiento identificar $u=\varphi(x)$ cuya derivada aparezca en el integrando (\cref{teo:sustitucion}).
\pasos
\begin{enumerate}[label=(\alph*)]
  \item $u=\ln x$, $\dif u=\frac{\dif x}{x}$: $\int u\dif u=\frac{(\ln x)^{2}}{2}+C$.
  \item $u=x^{2}-1$, $\dif u=2x\dif x$: $\frac12\int u^{1/2}\dif u=\frac13(x^{2}-1)^{3/2}+C$.
  \item $u=-x^{2}$, $\dif u=-2x\dif x$: $-\frac12\int e^{u}\dif u=-\frac12e^{-x^{2}}+C$.
  \item $u=\cos4x$, $\dif u=-4\sin4x\dif x$: $-\frac14\int\frac{\dif u}{u}=-\frac14\ln\abs{\cos4x}+C$.
  \item $u=x^{3}$, $\dif u=3x^{2}\dif x$: $\frac13\int\frac{\dif u}{1+u^{2}}=\frac13\arctan(x^{3})+C$.
\end{enumerate}
\resultado las primitivas anteriores.
\comprobacion por ejemplo, en (b), $\frac{d}{dx}\frac13(x^{2}-1)^{3/2}=\frac13\cdot\frac32(x^{2}-1)^{1/2}\cdot2x=x\sqrt{x^{2}-1}$; en (e), $\frac13\cdot\frac{3x^{2}}{1+x^{6}}$.
\end{solucion}

\begin{solucion}{eje:3-limites}
\planteamiento cambiar la variable y los límites a la vez (\cref{teo:sustitucion}).
\pasos
\begin{enumerate}[label=(\alph*)]
  \item $u=\ln x$: $x=1\mapsto u=0$, $x=e\mapsto u=1$. $\int_0^{1}\sqrt u\dif u=\eval{\frac23u^{3/2}}{0}{1}=\frac23$.
  \item $u=2+\cos x$, $\dif u=-\sin x\dif x$: $x=0\mapsto u=3$, $x=\pi\mapsto u=1$. $\int_3^{1}-u^{-3}\dif u=\int_1^{3}u^{-3}\dif u=\eval{-\frac{1}{2u^{2}}}{1}{3}=-\frac{1}{18}+\frac12=\frac49$.
  \item Para $0<b<1$, $x=\sin u$ con $u\in[0,\arcsin b]$, $\dif x=\cos u\dif u$ y $\sqrt{1-x^{2}}=\cos u>0$: $\int_0^{b}\frac{\dif x}{\sqrt{1-x^{2}}}=\int_0^{\arcsin b}\dif u=\arcsin b$, que tiende a $\frac\pi2$ cuando $b\to1^{-}$.
\end{enumerate}
\resultado (a) $\frac23$; (b) $\frac49$; (c) $\frac\pi2$.
\comprobacion (c) coincide con $\eval{\arcsin x}{0}{1}=\frac\pi2$; (b) es positiva, como corresponde a un integrando positivo en $(0,\pi)$.
\end{solucion}

\begin{solucion}{eje:3-partes}
\planteamiento \cref{teo:partes}; $u$ es el factor que se simplifica al derivar.
\pasos
\begin{enumerate}[label=(\alph*)]
  \item $u=x$, $v=-\cos x$: $-x\cos x+\int\cos x\dif x=-x\cos x+\sin x+C$.
  \item $u=x$, $v=\frac{e^{4x}}{4}$: $\frac{xe^{4x}}{4}-\frac{e^{4x}}{16}+C=\frac{(4x-1)e^{4x}}{16}+C$.
  \item $u=\ln x$, $v=\frac{x^{2}}{2}$: $\frac{x^{2}}{2}\ln x-\int\frac x2\dif x=\frac{x^{2}\ln x}{2}-\frac{x^{2}}{4}+C$.
  \item $u=\arctan x$, $v=x$: $x\arctan x-\int\frac{x}{1+x^{2}}\dif x=x\arctan x-\frac12\ln(1+x^{2})+C$.
  \item $u=x^{2}$, $v=e^{x}$: $x^{2}e^{x}-2\int xe^{x}\dif x=x^{2}e^{x}-2(x-1)e^{x}=(x^{2}-2x+2)e^{x}+C$, usando el \cref{ej:partes}(a).
\end{enumerate}
\resultado las primitivas anteriores.
\comprobacion derivando (e): $(2x-2)e^{x}+(x^{2}-2x+2)e^{x}=x^{2}e^{x}$; derivando (d): $\arctan x+\frac{x}{1+x^{2}}-\frac{x}{1+x^{2}}=\arctan x$.
\end{solucion}

\begin{solucion}{eje:3-ciclicas}
\planteamiento integrar dos veces por partes hasta que reaparezca la integral, y despejarla.
\pasos
\begin{enumerate}[label=\arabic*.]
  \item (a) Sea $K=\int e^{x}\cos x\dif x$. Con $u=\cos x$, $v=e^{x}$: $K=e^{x}\cos x+\int e^{x}\sin x\dif x$.
  \item Con $u=\sin x$, $v=e^{x}$: $\int e^{x}\sin x\dif x=e^{x}\sin x-K$. Luego $K=e^{x}\cos x+e^{x}\sin x-K$ y $K=\frac12e^{x}(\sin x+\cos x)+C$.
  \item (b) Sea $M=\int\cos(\ln x)\dif x$, $x>0$. Con $u=\cos(\ln x)$, $v=x$: $M=x\cos(\ln x)+\int\sin(\ln x)\dif x$.
  \item Con $u=\sin(\ln x)$, $v=x$: $\int\sin(\ln x)\dif x=x\sin(\ln x)-M$. Luego $M=\frac x2\bigl(\cos(\ln x)+\sin(\ln x)\bigr)+C$.
\end{enumerate}
\resultado (a) $\frac{e^{x}}{2}(\sin x+\cos x)+C$; (b) $\frac x2\bigl(\cos\ln x+\sin\ln x\bigr)+C$.
\comprobacion derivando (b): $\frac12(\cos\ln x+\sin\ln x)+\frac x2\bigl(-\frac{\sin\ln x}{x}+\frac{\cos\ln x}{x}\bigr)=\cos(\ln x)$. El cambio $x=e^{t}$ transforma (b) en $\int e^{t}\cos t\dif t$, que es (a).
\end{solucion}

\begin{solucion}{eje:3-fracciones}
\planteamiento \cref{prop:raices-simples} para raíces simples; coeficientes indeterminados para la raíz doble.
\pasos
\begin{enumerate}[label=(\alph*)]
  \item $A_3=\frac{1}{3-2}=1$, $A_2=\frac{1}{2-3}=-1$: $\int=\ln\abs{x-3}-\ln\abs{x-2}+C=\ln\Bigl|\frac{x-3}{x-2}\Bigr|+C$.
  \item $A_3=\frac{3}{1}=3$, $A_2=\frac{2}{-1}=-2$: $\int=3\ln\abs{x-3}-2\ln\abs{x-2}+C$.
  \item $x^{3}-x=x(x-1)(x+1)$: $A_0=\frac{1}{(-1)(1)}=-1$, $A_1=\frac{1}{1\cdot2}=\frac12$, $A_{-1}=\frac{1}{(-1)(-2)}=\frac12$: $\int=-\ln\abs x+\frac12\ln\abs{x-1}+\frac12\ln\abs{x+1}+C=\frac12\ln\abs{x^{2}-1}-\ln\abs x+C$.
  \item $\frac{3x+1}{(x-1)^{2}}=\frac{A}{x-1}+\frac{B}{(x-1)^{2}}$, es decir, $3x+1=A(x-1)+B$. En $x=1$: $B=4$; coeficiente de $x$: $A=3$. $\int=3\ln\abs{x-1}-\frac{4}{x-1}+C$.
\end{enumerate}
\resultado las primitivas anteriores.
\comprobacion recombinando (b): $\frac{3}{x-3}-\frac{2}{x-2}=\frac{3(x-2)-2(x-3)}{(x-3)(x-2)}=\frac{x}{(x-3)(x-2)}$.
\end{solucion}

\begin{solucion}{eje:3-fracciones-cuad}
\planteamiento \cref{teo:fracciones} con un factor cuadrático irreducible $x^{2}+1$.
\pasos
\begin{enumerate}[label=(\alph*)]
  \item $\frac{1}{(x-1)(x^{2}+1)}=\frac{A}{x-1}+\frac{Bx+C}{x^{2}+1}$, es decir, $1=A(x^{2}+1)+(Bx+C)(x-1)$. En $x=1$: $A=\frac12$. Coeficiente de $x^{2}$: $A+B=0$, $B=-\frac12$. Término independiente: $A-C=1$, $C=-\frac12$. Entonces
  $\int=\frac12\ln\abs{x-1}-\frac14\ln(x^{2}+1)-\frac12\arctan x+C$.
  \item $x^{4}-1=(x^{2}-1)(x^{2}+1)$ y $\frac{1}{x^{2}-1}-\frac{1}{x^{2}+1}=\frac{2}{x^{4}-1}$, luego $\frac{1}{x^{4}-1}=\frac12\cdot\frac{1}{x^{2}-1}-\frac12\cdot\frac{1}{x^{2}+1}$. Con el \cref{ej:fracciones}(a): $\int=\frac14\ln\Bigl|\frac{x-1}{x+1}\Bigr|-\frac12\arctan x+C$.
\end{enumerate}
\resultado las primitivas anteriores.
\comprobacion en (a), con $x=0$: $\frac{1}{(0-1)(0+1)}=-1$ y $\frac{1/2}{-1}+\frac{-1/2}{1}=-1$.
\end{solucion}

\begin{solucion}{eje:3-racional-impropia}
\planteamiento (a) dividir y separar la derivada del denominador; (b) completar el cuadrado.
\pasos
\begin{enumerate}[label=\arabic*.]
  \item (a) $\frac{x^{2}+1}{x^{2}+2x+2}=1-\frac{2x+1}{x^{2}+2x+2}$ y $2x+1=(2x+2)-1$, luego el integrando es $1-\frac{2x+2}{x^{2}+2x+2}+\frac{1}{(x+1)^{2}+1}$.
  \item Integrando cada término: $x-\ln(x^{2}+2x+2)+\arctan(x+1)+C$ (el denominador es siempre positivo).
  \item (b) $x^{2}+4x+13=(x+2)^{2}+9$, y con $h=2$, $s=3$: $\int\frac{\dif x}{(x+2)^{2}+9}=\frac13\arctan\frac{x+2}{3}+C$.
\end{enumerate}
\resultado (a) $x-\ln(x^{2}+2x+2)+\arctan(x+1)+C$; (b) $\frac13\arctan\frac{x+2}{3}+C$.
\comprobacion derivando (a): $1-\frac{2x+2}{x^{2}+2x+2}+\frac{1}{x^{2}+2x+2}=\frac{x^{2}+2x+2-2x-2+1}{x^{2}+2x+2}=\frac{x^{2}+1}{x^{2}+2x+2}$.
\end{solucion}

\begin{solucion}{eje:3-trig}
\planteamiento potencia impar: cambio de variable; potencias pares: fórmulas del ángulo doble.
\pasos
\begin{enumerate}[label=(\alph*)]
  \item $\sin^{3}x=(1-\cos^{2}x)\sin x$; con $u=\cos x$: $-\int(1-u^{2})\dif u=-\cos x+\frac{\cos^{3}x}{3}+C$.
  \item $\int_0^{\pi}\frac{1-\cos2x}{2}\dif x=\frac\pi2-\eval{\frac{\sin2x}{4}}{0}{\pi}=\frac\pi2$.
  \item $\cos^{4}x=\Bigl(\frac{1+\cos2x}{2}\Bigr)^{2}=\frac14\Bigl(1+2\cos2x+\frac{1+\cos4x}{2}\Bigr)=\frac38+\frac{\cos2x}{2}+\frac{\cos4x}{8}$. Luego $\int\cos^{4}x\dif x=\frac{3x}{8}+\frac{\sin2x}{4}+\frac{\sin4x}{32}+C$, y $\int_0^{\pi/2}\cos^{4}x\dif x=\frac{3\pi}{16}$.
\end{enumerate}
\resultado (a) $-\cos x+\frac13\cos^{3}x+C$; (b) $\frac\pi2$; (c) $\frac{3x}{8}+\frac{\sin2x}{4}+\frac{\sin4x}{32}+C$ y $\frac{3\pi}{16}$.
\comprobacion (c) coincide con $I_4=\frac{3\pi}{16}$ del \cref{eje:3-wallis}, pues $\int_0^{\pi/2}\cos^{n}=\int_0^{\pi/2}\sin^{n}$ (cambio $x\mapsto\frac\pi2-x$).
\end{solucion}

\begin{solucion}{eje:3-trig-sust}
\planteamiento $\sqrt{a^{2}-x^{2}}$: $x=a\sin t$; $1+x^{2}$: $x=\tan t$.
\pasos
\begin{enumerate}[label=(\alph*)]
  \item $x=a\sin t$, $t\in[-\frac\pi2,\frac\pi2]$: $\int_{-\pi/2}^{\pi/2}a\cos t\cdot a\cos t\dif t=a^{2}\int_{-\pi/2}^{\pi/2}\cos^{2}t\dif t=\frac{\pi a^{2}}{2}$.
  \item $x=\sin t$, $t\in(-\frac\pi2,\frac\pi2)$: $\int\frac{\sin^{2}t}{\cos t}\cos t\dif t=\int\sin^{2}t\dif t=\frac t2-\frac{\sin t\cos t}{2}+C=\frac12\bigl(\arcsin x-x\sqrt{1-x^{2}}\bigr)+C$.
  \item $x=\tan t$, $\dif x=(1+\tan^{2}t)\dif t$: $\int\frac{\dif t}{1+\tan^{2}t}=\int\cos^{2}t\dif t=\frac t2+\frac{\sin t\cos t}{2}+C$. Como $\sin t\cos t=\frac{\tan t}{1+\tan^{2}t}=\frac{x}{1+x^{2}}$, la primitiva es $\frac12\arctan x+\frac{x}{2(1+x^{2})}+C$.
\end{enumerate}
\resultado (a) $\frac{\pi a^{2}}{2}$; (b) $\frac12(\arcsin x-x\sqrt{1-x^{2}})+C$; (c) $\frac12\arctan x+\frac{x}{2(1+x^{2})}+C$.
\comprobacion (a) es el área de un semicírculo de radio $a$; con $a=1$ se recupera el \cref{ej:circulo}.
\end{solucion}

\begin{solucion}{eje:3-curso}
\datos (a) $a\neq0$, $\Delta>0$; (b) $n\ge2$.
\planteamiento (a) cambio lineal; (b) partes con $u=\ln x$ y paso al límite en el extremo $0$.
\pasos
\begin{enumerate}[label=\arabic*.]
  \item (a) $s=a(\Delta-u)$, $\dif s=-a\dif u$; $u=0\mapsto s=a\Delta$, $u=\Delta\mapsto s=0$: $\int_{a\Delta}^{0}e^{s}\bigl(-\frac{\dif s}{a}\bigr)=\frac1a\int_0^{a\Delta}e^{s}\dif s=\frac{e^{a\Delta}-1}{a}$.
  \item (b) Para $0<c<1$, con $u=\ln x$ y $v=\frac{x^{n}}{n}$: $\int_c^{1}x^{n-1}\ln x\dif x=\eval{\frac{x^{n}\ln x}{n}}{c}{1}-\frac1n\int_c^{1}x^{n-1}\dif x=-\frac{c^{n}\ln c}{n}-\frac{1-c^{n}}{n^{2}}$.
  \item Cuando $c\to0^{+}$, $c^{n}\ln c\to0$ y $c^{n}\to0$; el valor es $-\frac{1}{n^{2}}$.
\end{enumerate}
\resultado (a) $\frac{e^{a\Delta}-1}{a}$; (b) $-\frac{1}{n^{2}}$.
\comprobacion (a) tiende a $\Delta$ cuando $a\to0$, el valor con integrando constante~1. En (b), con $n=2$: $\int_0^{1}x\ln x\dif x=-\frac14$, negativo porque $\ln x<0$ en $(0,1)$.
\end{solucion}

\begin{solucion}{eje:3-reduccion}
\datos $a\neq0$, $n\in\N$; en (b), $b>0$.
\planteamiento (a) partes con $u=x^{n}$; (b) inducción usando (a) con $a=-1$ en $[0,b]$.
\pasos
\begin{enumerate}[label=\arabic*.]
  \item (a) Con $u=x^{n}$, $v=\frac{e^{ax}}{a}$: $\int x^{n}e^{ax}\dif x=\frac{x^{n}e^{ax}}{a}-\frac na\int x^{n-1}e^{ax}\dif x$.
  \item Con $a=1$: $J_0=e^{x}$, $J_1=xe^{x}-J_0=(x-1)e^{x}$, $J_2=x^{2}e^{x}-2J_1=(x^{2}-2x+2)e^{x}$, $J_3=x^{3}e^{x}-3J_2=(x^{3}-3x^{2}+6x-6)e^{x}+C$.
  \item (b) Sea $A_n=\int_0^{b}x^{n}e^{-x}\dif x$ y sea $P(n)$ la igualdad del enunciado. \emph{Caso base:} $A_0=1-e^{-b}=0!\,(1-e^{-b}\cdot1)$.
  \item \emph{Paso:} supóngase $P(n)$. Por (a) con $a=-1$ en forma definida, $A_{n+1}=\eval{-x^{n+1}e^{-x}}{0}{b}+(n+1)A_n=-b^{n+1}e^{-b}+(n+1)!\Bigl(1-e^{-b}\sum_{k=0}^{n}\frac{b^{k}}{k!}\Bigr)$.
  \item Como $b^{n+1}e^{-b}=(n+1)!\,e^{-b}\frac{b^{n+1}}{(n+1)!}$, se obtiene $A_{n+1}=(n+1)!\Bigl(1-e^{-b}\sum_{k=0}^{n+1}\frac{b^{k}}{k!}\Bigr)$, que es $P(n+1)$.
\end{enumerate}
\resultado $\int x^{3}e^{x}\dif x=(x^{3}-3x^{2}+6x-6)e^{x}+C$ y la fórmula de (b).
\comprobacion con $n=3$ y el intervalo $[0,1]$, (a) da $\int_0^{1}x^{3}e^{x}\dif x=(1-3+6-6)e+6=6-2e\approx0{,}5634$. En (b), cuando $b\to\infty$, $e^{-b}\sum_{k\le n}\frac{b^{k}}{k!}\to0$ y se obtiene $\int_0^{\infty}x^{n}e^{-x}\dif x=n!$, la \cref{prop:gamma-factorial}.
\end{solucion}

\begin{solucion}{eje:3-wallis}
\datos $I_n=\int_0^{\pi/2}\sin^{n}x\dif x$; $I_0=\frac\pi2$, $I_1=1$.
\planteamiento partes con $u=\sin^{n-1}x$ y $v'=\sin x$; después, inducción.
\pasos
\begin{enumerate}[label=\arabic*.]
  \item (a) Para $n\ge2$, con $v=-\cos x$: $I_n=\eval{-\sin^{n-1}x\cos x}{0}{\pi/2}+(n-1)\int_0^{\pi/2}\sin^{n-2}x\cos^{2}x\dif x$. El término de borde es nulo.
  \item Con $\cos^{2}x=1-\sin^{2}x$: $I_n=(n-1)(I_{n-2}-I_n)$, luego $nI_n=(n-1)I_{n-2}$.
  \item (b) Aplicando (a) $m$ veces hasta $I_0$ o $I_1$ se obtienen las fórmulas del enunciado (formalmente, por inducción sobre $m$). $I_4=\frac34\cdot\frac12\cdot\frac\pi2=\frac{3\pi}{16}$ e $I_5=\frac45\cdot\frac23\cdot1=\frac{8}{15}$.
  \item (c) Sea $Q(n)$: $nI_nI_{n-1}=\frac\pi2$. $Q(1)$: $1\cdot I_1I_0=\frac\pi2$. Si vale $Q(n)$, por (a), $(n+1)I_{n+1}I_n=(n+1)\cdot\frac{n}{n+1}I_{n-1}\cdot I_n=nI_nI_{n-1}=\frac\pi2$.
  \item En $[0,\frac\pi2]$, $0\le\sin x\le1$, luego $\sin^{n+1}x\le\sin^{n}x$ y, por la monotonía, $I_{n+1}\le I_n$.
\end{enumerate}
\resultado $I_n=\frac{n-1}{n}I_{n-2}$; $I_4=\frac{3\pi}{16}$, $I_5=\frac{8}{15}$; $nI_nI_{n-1}=\frac\pi2$ y $(I_n)$ decreciente.
\comprobacion $4\,I_4I_3=4\cdot\frac{3\pi}{16}\cdot\frac23=\frac\pi2$; además $I_5=0{,}533<I_4=0{,}589<I_3=0{,}667$.
\end{solucion}

\begin{solucion}{eje:3-beta}
\datos $B(m,n)=\int_0^{1}x^{m}(1-x)^{n}\dif x$, $m,n\in\N_0$.
\planteamiento inducción sobre $n$, para todos los $m$ a la vez; el paso integra por partes derivando $(1-x)^{n+1}$.
\pasos
\begin{enumerate}[label=\arabic*.]
  \item Sea $P(n)$: «$B(m,n)=\frac{m!\,n!}{(m+n+1)!}$ para todo $m\in\N_0$». \emph{Caso base:} $B(m,0)=\frac{1}{m+1}=\frac{m!\,0!}{(m+1)!}$.
  \item \emph{Paso:} supóngase $P(n)$ y sea $m\in\N_0$. Con $u=(1-x)^{n+1}$ y $v=\frac{x^{m+1}}{m+1}$:
  $B(m,n+1)=\eval{\frac{x^{m+1}(1-x)^{n+1}}{m+1}}{0}{1}+\frac{n+1}{m+1}\int_0^{1}x^{m+1}(1-x)^{n}\dif x$.
  \item El término de borde es nulo, y la integral es $B(m+1,n)$, a la que se aplica $P(n)$ con $m+1$:
  $B(m,n+1)=\frac{n+1}{m+1}\cdot\frac{(m+1)!\,n!}{(m+n+2)!}=\frac{m!\,(n+1)!}{(m+n+2)!}$, que es $P(n+1)$.
  \item Comprobación pedida: $x^{3}(1-x)^{2}=x^{3}-2x^{4}+x^{5}$, cuya integral es $\frac14-\frac25+\frac16=\frac{15-24+10}{60}=\frac{1}{60}$; la fórmula da $\frac{3!\,2!}{6!}=\frac{12}{720}=\frac{1}{60}$.
\end{enumerate}
\resultado $\int_0^{1}x^{m}(1-x)^{n}\dif x=\frac{m!\,n!}{(m+n+1)!}$.
\comprobacion la fórmula es simétrica en $m$ y $n$, como debe ser por el cambio $x\mapsto1-x$.
\end{solucion}

\begin{solucion}{eje:3-sustitucion-falsa}
\datos $I=\int_{-1}^{1}\frac{\dif x}{1+x^{2}}$; cambio $x=\varphi(t)=1/t$.
\planteamiento hacer el cálculo formal, comparar con el valor correcto y revisar las hipótesis del \cref{teo:sustitucion}.
\pasos
\begin{enumerate}[label=\arabic*.]
  \item (a) $\dif x=-\frac{\dif t}{t^{2}}$ y $\frac{1}{1+1/t^{2}}=\frac{t^{2}}{t^{2}+1}$, luego el integrando pasa a $-\frac{\dif t}{1+t^{2}}$. Los límites $x=\pm1$ corresponden a $t=\pm1$. Formalmente, $I=-\int_{-1}^{1}\frac{\dif t}{1+t^{2}}=-I$, y sería $I=0$.
  \item (b) $I=\arctan1-\arctan(-1)=\frac\pi2$.
  \item (c) El teorema exige que $\varphi$ tenga derivada continua en todo el intervalo $[\alpha,\beta]=[-1,1]$ de la variable $t$; pero $\varphi(t)=1/t$ no está definida en $t=0$. En la demostración se aplica la regla de Barrow a $F\circ\varphi$ en $[\alpha,\beta]$, que no es ni siquiera continua en $0$; ese es el paso que falla.
  \item En $\int_1^{2}\frac{\dif x}{1+x^{2}}$ se toma $t\in[\frac12,1]$, donde $\varphi(t)=1/t$ tiene derivada continua: $\int_1^{2}\frac{\dif x}{1+x^{2}}=\int_1^{1/2}\frac{-\dif t}{1+t^{2}}=\int_{1/2}^{1}\frac{\dif t}{1+t^{2}}=\frac\pi4-\arctan\frac12$.
\end{enumerate}
\resultado $I=\frac\pi2$; el cambio no es válido porque $\varphi$ no es de clase $C^{1}$ en $[-1,1]$; en $[1,2]$ sí lo es.
\comprobacion directamente, $\int_1^{2}\frac{\dif x}{1+x^{2}}=\arctan2-\frac\pi4\approx0{,}3218$, y $\frac\pi4-\arctan\frac12\approx0{,}7854-0{,}4636=0{,}3218$; ambos coinciden porque $\arctan2+\arctan\frac12=\frac\pi2$.
\end{solucion}

% ---------------------------------------------------------------------
\section{Tema 4}

\begin{solucion}{eje:4-clasificar}
\planteamiento buscar intervalos no acotados (primera especie) y puntos donde el integrando no está acotado (segunda especie).
\pasos
\begin{enumerate}[label=(\alph*)]
  \item Segunda especie en $x=0$: $\frac{1}{\sqrt x}\to\infty$.
  \item Segunda especie en el punto interior $x=1$; hay que estudiar $\int_0^{1}$ y $\int_1^{2}$ por separado.
  \item Primera especie (intervalo $[1,\infty)$); el integrando es acotado.
  \item Segunda especie en $x=\frac\pi2$: $\tan x\to\infty$.
  \item No es impropia: $\frac{\sin x}{x}\to1$ cuando $x\to0$, el integrando está acotado y, definiéndolo como~1 en $0$, es continuo en $[0,1]$.
  \item Primera especie en los dos extremos.
\end{enumerate}
\resultado impropias de segunda especie: (a), (b), (d); de primera: (c), (f); (e) no es impropia.
\comprobacion en (e), por el \cref{eje:1-un-punto}, el valor asignado en $x=0$ no afecta a la integral.
\end{solucion}

\begin{solucion}{eje:4-calcular}
\planteamiento \cref{prop:p-infinito,prop:p-cero} y cálculo directo del límite.
\pasos
\begin{enumerate}[label=(\alph*)]
  \item $p=2>1$: $\frac{1}{p-1}=1$.
  \item $p=10$: $\frac19$. Es el área bajo $y=x^{-10}$ a la derecha de $x=1$.
  \item $p=1$ en el origen: diverge ($-\ln c\to\infty$).
  \item $p=\frac12<1$: $\frac{1}{1-1/2}=2$.
  \item $\int_c^{e}\ln x\dif x=\eval{x\ln x-x}{c}{e}=(e-e)-(c\ln c-c)\to0$. Converge y vale~$0$.
  \item $\int_0^{b}xe^{-x^{2}}\dif x=\eval{-\frac12e^{-x^{2}}}{0}{b}=\frac12-\frac12e^{-b^{2}}\to\frac12$.
\end{enumerate}
\resultado (a) $1$; (b) $\frac19$; (c) diverge; (d) $2$; (e) $0$; (f) $\frac12$.
\comprobacion (e) vale $0$ porque el área negativa en $(0,1)$, que es $-1$ (\cref{ej:log-impropia}), se compensa con la positiva en $(1,e)$, que vale $\eval{x\ln x-x}{1}{e}=1$.
\end{solucion}

\begin{solucion}{eje:4-logaritmos}
\planteamiento cambios $u=\ln x$ y $u=\ln(\ln x)$, con $\dif u=\frac{\dif x}{x}$ y $\dif u=\frac{\dif x}{x\ln x}$.
\pasos
\begin{enumerate}[label=(\alph*)]
  \item $\int_2^{b}\frac{\dif x}{x\ln x}=\ln(\ln b)-\ln(\ln2)\to\infty$: diverge.
  \item $\int_2^{b}\frac{\dif x}{x(\ln x)^{2}}=\eval{-\frac{1}{\ln x}}{2}{b}=\frac{1}{\ln2}-\frac{1}{\ln b}\to\frac{1}{\ln2}\approx1{,}443$: converge.
  \item $\int_{100}^{b}\frac{\dif x}{x\ln x\ln(\ln x)}=\ln\bigl(\ln(\ln b)\bigr)-\ln\bigl(\ln(\ln100)\bigr)\to\infty$: diverge.
\end{enumerate}
\resultado (a) diverge; (b) converge a $\frac{1}{\ln2}$; (c) diverge.
\comprobacion la divergencia de (c) es lentísima: $\ln(\ln(\ln100))\approx0{,}42$, y para que la integral alcance~1 hace falta que $\ln(\ln(\ln b))\approx1{,}42$, es decir, $b\approx e^{e^{4{,}15}}$, un número de 28 cifras. Cada factor logarítmico más frena la divergencia, pero no la detiene; en cambio, un exponente $1+\varepsilon$ en el último logaritmo, como en (b), produce convergencia. No hay una frontera nítida entre ambos comportamientos.
\end{solucion}

\begin{solucion}{eje:4-comparar}
\planteamiento \cref{teo:comparacion-int,teo:comparacion-limite} con funciones $x^{-p}$.
\pasos
\begin{enumerate}[label=(\alph*)]
  \item $0<\frac{1}{x^{3}+1}\le\frac{1}{x^{3}}$ y $p=3>1$: \textbf{converge}.
  \item $\frac{2+\sin x}{x}\ge\frac1x\ge0$ y $\int_1^{\infty}\frac{\dif x}{x}$ diverge: \textbf{diverge}.
  \item $x+x^{3}\ge x$, luego $0<\frac{1}{\sqrt{x+x^{3}}}\le\frac{1}{\sqrt x}$, y $\int_0^{1}x^{-1/2}$ converge: \textbf{converge}.
  \item $\frac{1/\sqrt{x^{2}+x}}{1/x}=\frac{1}{\sqrt{1+1/x}}\to1$ y $\int_1^{\infty}\frac{\dif x}{x}$ diverge: \textbf{diverge}.
  \item $\frac{e^{x}}{x}\ge\frac1x>0$ en $(0,1]$ y $\int_0^{1}\frac{\dif x}{x}$ diverge: \textbf{diverge}.
  \item $x^{4}e^{-x}\to0$, luego $x^{2}e^{-x}\le\frac{1}{x^{2}}$ para $x\ge x_0$; en $[0,x_0]$ la integral es propia: \textbf{converge}.
\end{enumerate}
\resultado convergen (a), (c) y (f); divergen (b), (d) y (e).
\comprobacion (f) es $\Gamma(3)=2!=2$ (\cref{prop:gamma-factorial}).
\end{solucion}

\begin{solucion}{eje:4-seno}
\datos $\int_1^{N}\frac{\sin x}{x^{2}}\dif x$, $N>1$.
\planteamiento acotar el valor absoluto del integrando por $\frac{1}{x^{2}}$.
\pasos
\begin{enumerate}[label=\arabic*.]
  \item Por el \cref{teo:propiedades}(v) y la monotonía: $\int_1^{N}\frac{\sin x}{x^{2}}\dif x\le\int_1^{N}\frac{\abs{\sin x}}{x^{2}}\dif x\le\int_1^{N}\frac{\dif x}{x^{2}}=1-\frac1N<1$.
  \item $\int_1^{\infty}\frac{\abs{\sin x}}{x^{2}}$ converge por comparación con $\int_1^{\infty}x^{-2}$, y el \cref{teo:absoluta-int} da la convergencia de $\int_1^{\infty}\frac{\sin x}{x^{2}}$.
\end{enumerate}
\resultado la cota es estricta para todo $N$ y la integral impropia converge (absolutamente), con valor de módulo menor o igual que~1.
\comprobacion numéricamente, $\int_1^{\infty}\frac{\sin x}{x^{2}}\dif x\approx0{,}506$, dentro de la cota.
\end{solucion}

\begin{solucion}{eje:4-gamma-valores}
\planteamiento \cref{prop:gamma-factorial} y cambio lineal.
\pasos
\begin{enumerate}[label=\arabic*.]
  \item $\Gamma(2)=1!=1$; $\Gamma(5)=4!=24$; $\frac{\Gamma(7)}{\Gamma(5)}=\frac{6!}{4!}=6\cdot5=30$.
  \item Con $u=2x$, $x=\frac u2$, $\dif x=\frac{\dif u}{2}$: $\int_0^{\infty}x^{3}e^{-2x}\dif x=\int_0^{\infty}\frac{u^{3}}{8}e^{-u}\frac{\dif u}{2}=\frac{\Gamma(4)}{16}=\frac{6}{16}=\frac38$.
  \item En general, con $u=\lambda x$: $\int_0^{\infty}x^{n}e^{-\lambda x}\dif x=\frac{1}{\lambda^{n+1}}\int_0^{\infty}u^{n}e^{-u}\dif u=\frac{n!}{\lambda^{n+1}}$.
\end{enumerate}
\resultado $1$, $24$, $30$; $\frac38$; $\frac{n!}{\lambda^{n+1}}$.
\comprobacion con $n=0$: $\int_0^{\infty}e^{-\lambda x}\dif x=\frac1\lambda$, el \cref{ej:impropias1}(a).
\end{solucion}

\begin{solucion}{eje:4-gamma-medio}
\planteamiento cambio $x=u^{2}$ en la definición de $\Gamma(\frac12)$ y recurrencia de la \cref{prop:gamma-factorial}.
\pasos
\begin{enumerate}[label=\arabic*.]
  \item (a) En $[c,b]\subset(0,\infty)$, con $x=u^{2}$, $\dif x=2u\dif u$: $\int_c^{b}x^{-1/2}e^{-x}\dif x=\int_{\sqrt c}^{\sqrt b}\frac1ue^{-u^{2}}2u\dif u=2\int_{\sqrt c}^{\sqrt b}e^{-u^{2}}\dif u$.
  \item Pasando al límite ($c\to0^{+}$, $b\to\infty$) y usando el \cref{teo:gauss} y la paridad: $\Gamma(\frac12)=2\cdot\frac{\sqrt\pi}{2}=\sqrt\pi$.
  \item (b) $\Gamma(\frac32)=\frac12\Gamma(\frac12)=\frac{\sqrt\pi}{2}\approx0{,}8862$; $\Gamma(\frac52)=\frac32\Gamma(\frac32)=\frac{3\sqrt\pi}{4}\approx1{,}329$.
  \item En la notación $s!=\Gamma(s+1)$: $(\frac12)!=\Gamma(\frac32)=\frac{\sqrt\pi}{2}$.
\end{enumerate}
\resultado $\Gamma(\frac12)=\sqrt\pi$, $\Gamma(\frac32)=\frac{\sqrt\pi}{2}$, $\Gamma(\frac52)=\frac{3\sqrt\pi}{4}$.
\comprobacion $\Gamma(\frac32)\approx0{,}886$ está entre $\Gamma(1)=1$ y $\Gamma(2)=1$, cerca del mínimo de $\Gamma$ en $(1,2)$, que vale aproximadamente $0{,}886$.
\end{solucion}

\begin{solucion}{eje:4-erf}
\datos $E(x)=\frac{1}{\sqrt{2\pi}}\int_0^{x}e^{-u^{2}/2}\dif u$, $F(x)=\int_0^{x}e^{-t^{2}}\dif t$, $F(x)\to\frac{\sqrt\pi}{2}$.
\planteamiento cambio $u=\sqrt2\,t$.
\pasos
\begin{enumerate}[label=\arabic*.]
  \item Con $u=\sqrt2\,t$, $\dif u=\sqrt2\dif t$, el límite $x$ pasa a $x/\sqrt2$: $E(x)=\frac{\sqrt2}{\sqrt{2\pi}}\int_0^{x/\sqrt2}e^{-t^{2}}\dif t=\frac{1}{\sqrt\pi}F\bigl(\tfrac{x}{\sqrt2}\bigr)$.
  \item Cuando $x\to\infty$: $E(x)\to\frac{1}{\sqrt\pi}\cdot\frac{\sqrt\pi}{2}=\frac12$.
  \item El integrando $e^{-t^{2}}$ es par, luego $F$ es impar (\cref{eje:2-par-impar}) y $E$ también: $E(x)\to-\frac12$ cuando $x\to-\infty$.
  \item (b) Por la aditividad, $\frac{1}{\sqrt{2\pi}}\int_a^{b}e^{-u^{2}/2}\dif u=E(b)-E(a)$.
\end{enumerate}
\resultado $E(x)=\frac{1}{\sqrt\pi}F(x/\sqrt2)$; límites $\pm\frac12$; $\frac{1}{\sqrt{2\pi}}\int_a^{b}e^{-u^{2}/2}\dif u=E(b)-E(a)$.
\comprobacion $\frac{1}{\sqrt{2\pi}}\int_{-N}^{N}e^{-u^{2}/2}\dif u=E(N)-E(-N)=2E(N)\to1$: la densidad normal estándar integra~1 (\cref{cor:gauss}). La función de distribución es $\Phi(x)=\frac12+E(x)$.
\end{solucion}

\begin{solucion}{eje:4-normal}
\planteamiento integración por partes en $[-R,R]$ o $[0,b]$ y paso al límite.
\pasos
\begin{enumerate}[label=\arabic*.]
  \item (a) Con $u=x$, $v=-e^{-x^{2}/2}$: $\int_{-R}^{R}x^{2}e^{-x^{2}/2}\dif x=\eval{-xe^{-x^{2}/2}}{-R}{R}+\int_{-R}^{R}e^{-x^{2}/2}\dif x=-2Re^{-R^{2}/2}+\int_{-R}^{R}e^{-x^{2}/2}\dif x$.
  \item El integrando es par y positivo; ambas colas convergen (por comparación, $x^{2}e^{-x^{2}/2}\le\frac{1}{x^{2}}$ para $x$ grande), de modo que basta el límite simétrico. Cuando $R\to\infty$: $Re^{-R^{2}/2}\to0$ y, por el \cref{cor:gauss}, la integral vale $\sqrt{2\pi}$.
  \item Para la normal estándar, $\E[X]=\frac{1}{\sqrt{2\pi}}\int xe^{-x^{2}/2}\dif x=0$ (integrando impar y absolutamente integrable) y $\E[X^{2}]=\frac{1}{\sqrt{2\pi}}\sqrt{2\pi}=1$; luego $\operatorname{Var}(X)=1$.
  \item (b) Con $u=x$, $v=-\frac12e^{-x^{2}}$: $\int_0^{b}x^{2}e^{-x^{2}}\dif x=-\frac b2e^{-b^{2}}+\frac12\int_0^{b}e^{-x^{2}}\dif x\to\frac12\cdot\frac{\sqrt\pi}{2}=\frac{\sqrt\pi}{4}$.
\end{enumerate}
\resultado $\int_{-\infty}^{\infty}x^{2}e^{-x^{2}/2}\dif x=\sqrt{2\pi}$, varianza~1; $\int_0^{\infty}x^{2}e^{-x^{2}}\dif x=\frac{\sqrt\pi}{4}\approx0{,}4431$.
\comprobacion (b) también se obtiene de (a) con el cambio $x=t/\sqrt2$: $\int_0^{\infty}x^{2}e^{-x^{2}}\dif x=\frac{1}{2\sqrt2}\int_0^{\infty}t^{2}e^{-t^{2}/2}\dif t=\frac{1}{2\sqrt2}\cdot\frac{\sqrt{2\pi}}{2}=\frac{\sqrt\pi}{4}$.
\end{solucion}

\begin{solucion}{eje:4-gamma-conv}
\datos $s\le0$.
\planteamiento acotar inferiormente el integrando en $(0,1]$ por una función $x^{-p}$ con $p\ge1$.
\pasos
\begin{enumerate}[label=\arabic*.]
  \item En $(0,1]$, $e^{-x}\ge e^{-1}$, luego $x^{s-1}e^{-x}\ge e^{-1}x^{-(1-s)}>0$.
  \item Como $s\le0$, el exponente $p=1-s\ge1$, y $\int_0^{1}x^{-p}\dif x$ diverge (\cref{prop:p-cero}).
  \item Por el \cref{teo:comparacion-int} (en su versión para el extremo $0$), $\int_0^{1}x^{s-1}e^{-x}\dif x$ diverge, y con ella la integral en $(0,\infty)$, cuyo integrando es positivo.
  \item Junto con la convergencia para $s>0$ probada tras la \cref{def:gamma}, la integral converge si y solo si $s>0$.
\end{enumerate}
\resultado la integral que define $\Gamma(s)$ converge exactamente para $s>0$.
\comprobacion con $s=0$ el integrando es $\frac{e^{-x}}{x}\approx\frac1x$ cerca de $0$, el caso límite de la divergencia logarítmica.
\end{solucion}

\begin{solucion}{eje:4-limite-comp}
\datos $f,g>0$ en $[a,\infty)$, $\frac fg\to L\in(0,\infty)$.
\planteamiento con $\varepsilon=\frac L2$ en la definición de límite se obtienen cotas $\frac L2g\le f\le2Lg$; después se aplica el \cref{teo:comparacion-int}.
\pasos
\begin{enumerate}[label=\arabic*.]
  \item Existe $x_0\ge a$ tal que $\bigl|\frac{f(x)}{g(x)}-L\bigr|<\frac L2$ para $x\ge x_0$, es decir, $\frac L2g(x)<f(x)<\frac{3L}{2}g(x)\le2Lg(x)$.
  \item Si $\int_a^{\infty}g$ converge, también $\int_{x_0}^{\infty}2Lg$, y por comparación $\int_{x_0}^{\infty}f$ converge; $\int_a^{x_0}f$ es una integral propia, luego $\int_a^{\infty}f$ converge.
  \item Si $\int_a^{\infty}f$ converge, de $g<\frac2Lf$ en $[x_0,\infty)$ se deduce del mismo modo que $\int_a^{\infty}g$ converge.
  \item Por tanto ambas convergen o ambas divergen.
  \item Aplicaciones: $\frac{(x+1)/(x^{3}+2)}{1/x^{2}}=\frac{x^{3}+x^{2}}{x^{3}+2}\to1$ y $\int_1^{\infty}x^{-2}$ converge: la primera \textbf{converge}. Para $x\ge1$, $\sin\frac1x>0$ y $\frac{\sin(1/x)}{1/x}\to1$ (pues $\frac{\sin u}{u}\to1$ cuando $u\to0$); como $\int_1^{\infty}\frac{\dif x}{x}$ diverge, la segunda \textbf{diverge}.
\end{enumerate}
\resultado el teorema queda demostrado; la primera integral converge y la segunda diverge.
\comprobacion la hipótesis $L>0$ se usa en el paso~3 (para dividir por $L$); con $L=0$ solo vale la implicación del paso~2.
\end{solucion}

\begin{solucion}{eje:4-dirichlet}
\datos $\int_1^{\infty}\frac{\sin x}{x}\dif x$.
\planteamiento integrar por partes para pasar de $\frac1x$ a $\frac{1}{x^{2}}$, que es integrable en el infinito.
\pasos
\begin{enumerate}[label=\arabic*.]
  \item Con $u=\frac1x$, $v=-\cos x$: $\int_1^{b}\frac{\sin x}{x}\dif x=\eval{-\frac{\cos x}{x}}{1}{b}-\int_1^{b}\frac{\cos x}{x^{2}}\dif x=\cos1-\frac{\cos b}{b}-\int_1^{b}\frac{\cos x}{x^{2}}\dif x$.
  \item $\bigl|\frac{\cos b}{b}\bigr|\le\frac1b\to0$.
  \item $\bigl|\frac{\cos x}{x^{2}}\bigr|\le\frac{1}{x^{2}}$, luego $\int_1^{\infty}\frac{\cos x}{x^{2}}\dif x$ converge absolutamente (como en el \cref{ej:criterios}(b)).
  \item Los tres términos tienen límite cuando $b\to\infty$, y $\int_1^{\infty}\frac{\sin x}{x}\dif x=\cos1-\int_1^{\infty}\frac{\cos x}{x^{2}}\dif x$ converge.
  \item En $(0,1]$, $\frac{\sin x}{x}$ está acotada y tiende a~1 en $0$; definida como~1 en $x=0$ es continua en $[0,1]$, y $\int_0^{1}$ es una integral ordinaria. Luego $\int_0^{\infty}\frac{\sin x}{x}\dif x$ converge.
\end{enumerate}
\resultado la integral converge.
\comprobacion la convergencia no es absoluta: en cada $[k\pi,(k+1)\pi]$, $\int\frac{\abs{\sin x}}{x}\ge\frac{2}{(k+1)\pi}$, y la suma de estas cotas es una serie armónica, divergente (\cref{teo:armonica}). Su valor, que se obtiene con otras técnicas, es $\frac\pi2$.
\end{solucion}

\begin{solucion}{eje:4-potencias}
\datos $f(x)=(\alpha-1)x^{-\alpha}$, $x\ge1$.
\planteamiento \cref{prop:p-infinito} con exponentes $\alpha$, $\alpha-1$ y $\alpha-2$.
\pasos
\begin{enumerate}[label=\arabic*.]
  \item (a) Si $\alpha<1$, $f<0$; si $\alpha=1$, $f\equiv0$ e integra $0$. Si $\alpha>1$, $f>0$ e $\int_1^{\infty}f=(\alpha-1)\cdot\frac{1}{\alpha-1}=1$. Luego $f$ es densidad si y solo si $\alpha>1$.
  \item (b) $\int_1^{\infty}xf(x)\dif x=(\alpha-1)\int_1^{\infty}x^{-(\alpha-1)}\dif x$ converge si y solo si $\alpha-1>1$, es decir, $\alpha>2$, y vale $\frac{\alpha-1}{\alpha-2}$.
  \item (c) $\int_1^{\infty}x^{2}f(x)\dif x=(\alpha-1)\int_1^{\infty}x^{-(\alpha-2)}\dif x$ es finito si y solo si $\alpha>3$, y vale $\frac{\alpha-1}{\alpha-3}$.
\end{enumerate}
\resultado densidad si $\alpha>1$; media finita si $\alpha>2$, igual a $\frac{\alpha-1}{\alpha-2}$; segundo momento finito si $\alpha>3$.
\comprobacion con $\alpha=3{,}5$, la media es $\frac{2{,}5}{1{,}5}=\frac53$. Con $\alpha=2{,}5$ la media es finita ($3$) pero la varianza es infinita: las colas pesadas de las leyes de potencia hacen que los promedios muestrales fluctúen mucho.
\end{solucion}

\begin{solucion}{eje:4-decreciente}
\datos $f\ge0$ decreciente en $[1,\infty)$, $\int_1^{\infty}f=I$ finita.
\planteamiento acotar $\frac x2f(x)$ por una integral sobre $[\frac x2,x]$, que tiende a~$0$.
\pasos
\begin{enumerate}[label=\arabic*.]
  \item Sea $F(y)=\int_1^{y}f$, con $F(y)\to I$. Para $x\ge2$, $\int_{x/2}^{x}f=F(x)-F(x/2)\to I-I=0$.
  \item Para $t\in[\frac x2,x]$, el decrecimiento da $f(t)\ge f(x)$; por la monotonía, $\int_{x/2}^{x}f\ge\frac x2f(x)$.
  \item Luego $0\le xf(x)\le2\bigl(F(x)-F(x/2)\bigr)\to0$, y $xf(x)\to0$.
  \item Con $f(x)=\frac1x$, positiva y decreciente, $xf(x)=1\not\to0$; por el contrarrecíproco, $\int_1^{\infty}\frac{\dif x}{x}$ no puede converger.
\end{enumerate}
\resultado $\lim_{x\to\infty}xf(x)=0$; en particular, $\int_1^{\infty}\frac{\dif x}{x}$ diverge.
\comprobacion la hipótesis de monotonía es necesaria: hay funciones continuas no negativas con integral convergente que no tienden a~$0$, formadas por picos triangulares de altura~1 y anchura cada vez menor.
\end{solucion}

\begin{solucion}{eje:4-wallis-gauss}
\datos $J=\int_0^{\infty}e^{-x^{2}}\dif x$; $I_n$ del \cref{eje:3-wallis}, con $nI_nI_{n-1}=\frac\pi2$ e $(I_n)$ decreciente.
\planteamiento encerrar $J$ entre dos integrales de Wallis y usar su comportamiento asintótico.
\pasos
\begin{enumerate}[label=\arabic*.]
  \item (a) De $e^{t}\ge1+t$ con $t=-x^{2}$: $e^{-x^{2}}\ge1-x^{2}$, y en $[0,1]$ ambos miembros son no negativos; elevando a $n$: $(1-x^{2})^{n}\le e^{-nx^{2}}$. Con $t=x^{2}$: $e^{x^{2}}\ge1+x^{2}$, luego $e^{-nx^{2}}\le(1+x^{2})^{-n}$.
  \item (b) Con el cambio $x=u/\sqrt n$, $\int_0^{\infty}e^{-nx^{2}}\dif x=\frac{J}{\sqrt n}$. Integrando las desigualdades de (a), y usando que el integrando es positivo:
  $\int_0^{1}(1-x^{2})^{n}\dif x\le\int_0^{1}e^{-nx^{2}}\dif x\le\frac{J}{\sqrt n}\le\int_0^{\infty}\frac{\dif x}{(1+x^{2})^{n}}$.
  \item Con $x=\cos\theta$, $\theta\in[0,\frac\pi2]$: $\int_0^{1}(1-x^{2})^{n}\dif x=\int_0^{\pi/2}\sin^{2n}\theta\,\sin\theta\dif\theta=I_{2n+1}$.
  \item Con $x=\tan\theta$, $\dif x=\frac{\dif\theta}{\cos^{2}\theta}$ y $(1+\tan^{2}\theta)^{-n}=\cos^{2n}\theta$: $\int_0^{\infty}\frac{\dif x}{(1+x^{2})^{n}}=\int_0^{\pi/2}\cos^{2n-2}\theta\dif\theta=I_{2n-2}$, pues $\int_0^{\pi/2}\cos^{m}=\int_0^{\pi/2}\sin^{m}$. Multiplicando por $\sqrt n$: $\sqrt n\,I_{2n+1}\le J\le\sqrt n\,I_{2n-2}$.
  \item (c) Por el decrecimiento, $I_n^{2}\le I_nI_{n-1}=\frac{\pi}{2n}$ e $I_n^{2}\ge I_nI_{n+1}=\frac{\pi}{2(n+1)}$, de donde $\sqrt{\frac{\pi}{2(n+1)}}\le I_n\le\sqrt{\frac{\pi}{2n}}$.
  \item Entonces $\sqrt n\,I_{2n+1}\ge\sqrt{\frac{\pi n}{4(n+1)}}$ y, para $n\ge2$, $\sqrt n\,I_{2n-2}\le\sqrt{\frac{\pi n}{4(n-1)}}$. Ambas cotas tienden a $\frac{\sqrt\pi}{2}$ cuando $n\to\infty$.
  \item Como $J$ no depende de $n$ y está entre dos cantidades que tienden a $\frac{\sqrt\pi}{2}$, $J=\frac{\sqrt\pi}{2}$.
\end{enumerate}
\resultado $\int_0^{\infty}e^{-x^{2}}\dif x=\frac{\sqrt\pi}{2}$, y por paridad $\int_{-\infty}^{\infty}e^{-x^{2}}\dif x=\sqrt\pi$.
\comprobacion con $n=20$, las cotas del paso~4 valen $0{,}8700\le J\le0{,}9033$, y $\frac{\sqrt\pi}{2}\approx0{,}8862$.
\end{solucion}

% ---------------------------------------------------------------------
\section{Tema 5}

\begin{solucion}{eje:5-vf}
\planteamiento contrastar con la \cref{prop:condicion-necesaria}, el \cref{teo:armonica}, el \cref{teo:absoluta-series} y la observación tras el \cref{teo:raiz}.
\pasos
\begin{enumerate}[label=(\alph*)]
  \item \textbf{Falsa}: $\frac1n\to0$ y la serie armónica diverge.
  \item \textbf{Verdadera}: \cref{prop:condicion-necesaria}.
  \item \textbf{Falsa}. La sucesión converge a $\ln1=0$, pero la serie es telescópica: $\ln(1+\frac1k)=\ln(k+1)-\ln k$, luego $S_n=\ln(n+1)\to\infty$ y diverge.
  \item \textbf{Verdadera}: \cref{teo:absoluta-series}.
  \item \textbf{Falsa}: $\sum\frac{1}{n^{2}}$ tiene cociente $\frac{n^{2}}{(n+1)^{2}}\to1$ y converge.
\end{enumerate}
\resultado (a) F; (b) V; (c) F; (d) V; (e) F.
\comprobacion (c) es otro contraejemplo a (a), con términos que se comportan como $\frac1k$ (comparación en el límite).
\end{solucion}

\begin{solucion}{eje:5-limites}
\planteamiento álgebra de límites, conjugado, encaje y logaritmos.
\pasos
\begin{enumerate}[label=(\alph*)]
  \item Dividiendo por $n^{2}$: $\frac{3+1/n^{2}}{2-1/n}\to\frac32$.
  \item $\sqrt{n+1}-\sqrt n=\frac{1}{\sqrt{n+1}+\sqrt n}\to0$.
  \item Por la regla de L'Hôpital aplicada a $\frac{x}{2^{x}}$: $\frac{1}{2^{x}\ln2}\to0$; luego $\frac{n}{2^{n}}\to0$.
  \item $\bigl|\frac{\sin n}{n}\bigr|\le\frac1n\to0$; por encaje, el límite es~$0$.
  \item $\ln\bigl(1+\frac2n\bigr)^{n}=n\ln\bigl(1+\frac2n\bigr)=2\cdot\frac{\ln(1+h)}{h}$ con $h=\frac2n\to0$, que tiende a $2$. Por la continuidad de la exponencial, el límite es $e^{2}$.
\end{enumerate}
\resultado (a) $\frac32$; (b) $0$; (c) $0$; (d) $0$; (e) $e^{2}\approx7{,}389$.
\comprobacion (e) con $n=1000$: $(1{,}002)^{1000}\approx7{,}374$.
\end{solucion}

\begin{solucion}{eje:5-geometricas}
\planteamiento \cref{teo:geometrica}: suma $=\frac{\text{primer término}}{1-r}$ si $\abs r<1$.
\pasos
\begin{enumerate}[label=\arabic*.]
  \item (a) Primer término $\frac13$, $r=\frac13$: $\frac{1/3}{2/3}=\frac12$. Primer término $1$, $r=-\frac14$: $\frac{1}{5/4}=\frac45$.
  \item (b) Primer término $3\cdot\frac{4}{25}=\frac{12}{25}$, $r=-\frac25$: $\frac{12/25}{7/5}=\frac{12}{35}$.
  \item (c) $0{,}8787\ldots=\frac{87}{100}+\frac{87}{100^{2}}+\cdots=\frac{87/100}{1-1/100}=\frac{87}{99}=\frac{29}{33}$. Del mismo modo, $0{,}123123\ldots=\frac{123/1000}{1-1/1000}=\frac{123}{999}=\frac{41}{333}$.
\end{enumerate}
\resultado (a) $\frac12$ y $\frac45$; (b) $\frac{12}{35}$; (c) $\frac{29}{33}$ y $\frac{41}{333}$.
\comprobacion $29/33=0{,}878787\ldots$ y $41/333=0{,}123123\ldots$ al dividir.
\end{solucion}

\begin{solucion}{eje:5-pelota}
\datos altura inicial $10$~m; tras cada bote sube a $\frac35$ de la altura anterior.
\planteamiento la primera caída es de 10~m; después, cada altura $h_k=10\bigl(\frac35\bigr)^{k}$, $k\ge1$, se recorre dos veces (subida y bajada).
\pasos
\begin{enumerate}[label=\arabic*.]
  \item Distancia: $10+2\sum_{k\ge1}10\bigl(\frac35\bigr)^{k}=10+20\cdot\frac{3/5}{2/5}=10+30=40$~m.
  \item Tiempo: cada trayecto de altura $h$ dura $\sqrt{2h/g}$. Total: $\sqrt{\frac{20}{g}}\Bigl(1+2\sum_{k\ge1}\bigl(\sqrt{3/5}\bigr)^{k}\Bigr)$, serie geométrica de razón $\sqrt{0{,}6}\approx0{,}7746<1$.
  \item Con $g=9{,}8$~m/s$^{2}$: $\sum_{k\ge1}r^{k}=\frac{r}{1-r}\approx3{,}436$, y el tiempo total es $1{,}429\cdot(1+6{,}873)\approx11{,}2$~s.
\end{enumerate}
\resultado 40~m; el tiempo total es finito, unos $11{,}2$~s.
\comprobacion infinitos botes caben en un tiempo finito porque las duraciones decrecen geométricamente.
\end{solucion}

\begin{solucion}{eje:5-comparacion}
\planteamiento \cref{teo:comparacion-series,teo:comparacion-limite-series} con series $p$.
\pasos
\begin{enumerate}[label=(\alph*)]
  \item $0<\frac{1}{n^{2}+10}\le\frac{1}{n^{2}}$: \textbf{converge}.
  \item $\frac{1}{2n-1}\ge\frac{1}{2n}$ y $\sum\frac{1}{2n}$ diverge: \textbf{diverge}.
  \item $\frac{(n+1)/(n^{3}+2)}{1/n^{2}}=\frac{n^{3}+n^{2}}{n^{3}+2}\to1$: \textbf{converge}.
  \item $\frac{1/\sqrt{n^{2}+n}}{1/n}=\frac{1}{\sqrt{1+1/n}}\to1$: \textbf{diverge}.
  \item $0<\sin\frac{1}{n^{2}}\le\frac{1}{n^{2}}$ (pues $\sin u\le u$ para $u\ge0$): \textbf{converge}.
\end{enumerate}
\resultado convergen (a), (c), (e); divergen (b), (d).
\comprobacion en (e) también vale la comparación en el límite: $\frac{\sin(1/n^{2})}{1/n^{2}}\to1$.
\end{solucion}

\begin{solucion}{eje:5-cociente}
\planteamiento \cref{teo:cociente,teo:raiz}.
\pasos
\begin{enumerate}[label=(\alph*)]
  \item $\frac{a_{n+1}}{a_n}=\frac{3}{n+1}\to0$: \textbf{converge}.
  \item $\frac{a_{n+1}}{a_n}=\frac{(n+1)!\,n^{n}}{(n+1)^{n+1}n!}=\frac{n^{n}}{(n+1)^{n}}=\frac{1}{(1+1/n)^{n}}\to\frac1e<1$: \textbf{converge}.
  \item $\frac{a_{n+1}}{a_n}=\frac{2n^{2}}{(n+1)^{2}}\to2>1$: \textbf{diverge}.
  \item $\sqrt[n]{a_n}=\frac{n}{2n+1}\to\frac12<1$: \textbf{converge}.
  \item $\frac{a_{n+1}}{a_n}=\frac{n^{3}}{(n+1)^{3}}\to1$: el criterio no decide; la serie converge por ser una serie $p$ con $p=3$.
\end{enumerate}
\resultado convergen (a), (b), (d), (e); diverge (c).
\comprobacion la suma de (a) desde $n=0$ es $e^{3}$ (\cref{teo:exp}); en (c), $\frac{2^{n}}{n^{2}}\to\infty$, de acuerdo con la divergencia.
\end{solucion}

\begin{solucion}{eje:5-integral}
\planteamiento \cref{teo:criterio-integral}; comprobar antes que $f$ es positiva y decreciente (al menos desde algún punto, lo que basta).
\pasos
\begin{enumerate}[label=(\alph*)]
  \item $f(x)=e^{-x}$; $\int_1^{\infty}e^{-x}\dif x=e^{-1}$: \textbf{converge}. Como geométrica: $\sum_{n\ge1}e^{-n}=\frac{e^{-1}}{1-e^{-1}}=\frac{1}{e-1}\approx0{,}582$.
  \item $f(x)=xe^{-x}$, $f'(x)=(1-x)e^{-x}\le0$ para $x\ge1$; $\int_1^{\infty}xe^{-x}\dif x=\eval{-(x+1)e^{-x}}{1}{\infty}=\frac2e$: \textbf{converge}.
  \item $f(x)=\frac{\ln x}{x}$, $f'(x)=\frac{1-\ln x}{x^{2}}<0$ para $x>e$; $\int_3^{b}\frac{\ln x}{x}\dif x=\frac{(\ln b)^{2}-(\ln3)^{2}}{2}\to\infty$: \textbf{diverge}.
  \item $f(x)=\frac{1}{x(\ln x)^{2}}$, decreciente en $[2,\infty)$; $\int_2^{\infty}f=\frac{1}{\ln2}$ (\cref{eje:4-logaritmos}): \textbf{converge}.
\end{enumerate}
\resultado convergen (a), (b), (d); diverge (c).
\comprobacion (c) también se obtiene por comparación: $\frac{\ln n}{n}\ge\frac1n$ para $n\ge3$. La suma de (b) es $\frac{e}{(e-1)^{2}}\approx0{,}9207$ (\cref{eje:6-derivar-geom} con $x=e^{-1}$).
\end{solucion}

\begin{solucion}{eje:5-armonica}
\datos $\ln(n+1)\le H_n\le1+\ln n$; $\gamma\approx0{,}5772$.
\planteamiento despejar $n$ en las cotas.
\pasos
\begin{enumerate}[label=\arabic*.]
  \item Si $H_n>7$, entonces $1+\ln n>7$, es decir, $n>e^{6}\approx403{,}4$. Hacen falta al menos 404 términos, más de 400.
  \item Si $\ln(n+1)>7$, es decir, $n+1>e^{7}\approx1096{,}6$, entonces $H_n>7$. Basta $n=1096$; con $1097$ términos ya se supera con seguridad.
  \item (b) $H_n\approx\ln n+\gamma>7$ da $n>e^{7-\gamma}\approx615{,}7$; la estimación es $n=616$.
\end{enumerate}
\resultado entre 404 y 1096 términos según las cotas; la estimación con $\gamma$ es $616$.
\comprobacion calculando las sumas parciales, el primer $n$ con $H_n>7$ es exactamente $616$.
\end{solucion}

\begin{solucion}{eje:5-alternadas}
\planteamiento estudiar primero $\sum\abs{a_n}$; si diverge, aplicar el \cref{teo:leibniz} o la \cref{prop:condicion-necesaria}.
\pasos
\begin{enumerate}[label=(\alph*)]
  \item $\sum\frac{1}{n!}$ converge (cociente $\frac{1}{n+1}\to0$): \textbf{absolutamente}.
  \item $\sum\frac{3^{n}}{n!}$ converge (\cref{eje:5-cociente}(a)): \textbf{absolutamente}.
  \item $\frac{n^{2}}{1+n^{4}}\le\frac{1}{n^{2}}$: \textbf{absolutamente}.
  \item $\sum\frac{1}{\sqrt n}$ diverge ($p=\frac12$), pero $\frac{1}{\sqrt n}$ decrece a~$0$: \textbf{condicionalmente}.
  \item $\frac{n}{n+1}\to1$, luego $(-1)^{n}\frac{n}{n+1}\not\to0$: \textbf{diverge}.
\end{enumerate}
\resultado absolutas (a), (b), (c); condicional (d); divergente (e).
\comprobacion la suma de (a) es $1-\frac1e$, de acuerdo con $e^{-1}=\sum\frac{(-1)^{n}}{n!}$ (\cref{teo:exp} con $x=-1$).
\end{solucion}

\begin{solucion}{eje:5-error-alternada}
\planteamiento cota $\abs{S-S_n}\le b_{n+1}$ del \cref{teo:leibniz}.
\pasos
\begin{enumerate}[label=\arabic*.]
  \item (a) $\abs{S_5-\ln2}\le b_6=\frac16\approx0{,}1667$. Real: $S_5=1-\frac12+\frac13-\frac14+\frac15=\frac{47}{60}\approx0{,}7833$ y $S_5-\ln2\approx0{,}0902$.
  \item (b) Se necesita $b_{n+1}=\frac{1}{(n+1)^{2}}<10^{-3}$, es decir, $n+1>\sqrt{1000}\approx31{,}6$, o sea $n\ge31$.
\end{enumerate}
\resultado (a) error $\le\frac16$ (real $0{,}090$); (b) 31 términos.
\comprobacion en (b), $\frac{1}{32^{2}}=\frac{1}{1024}<10^{-3}$, mientras que $\frac{1}{31^{2}}=\frac{1}{961}>10^{-3}$. La cota es conservadora: el error real suele ser aproximadamente la mitad de $b_{n+1}$.
\end{solucion}

\begin{solucion}{eje:5-descuento}
\datos $\gamma=0{,}9$; $\beta\in\{0{,}9;\ 0{,}999\}$.
\planteamiento series geométricas (\cref{teo:geometrica}) y logaritmos.
\pasos
\begin{enumerate}[label=\arabic*.]
  \item Recompensa constante: $G=\sum_{k\ge0}0{,}9^{k}=\frac{1}{0{,}1}=10$.
  \item Si $R_{k+1}=1$ solo cuando $k+1$ es impar, es decir, $k$ par: $G=\sum_{j\ge0}\gamma^{2j}=\frac{1}{1-\gamma^{2}}=\frac{1}{0{,}19}\approx5{,}263$.
  \item $1-\beta^{k}\ge0{,}99\iff\beta^{k}\le0{,}01\iff k\ge\frac{\ln0{,}01}{\ln\beta}$ (al dividir por $\ln\beta<0$ cambia el sentido).
  \item $\beta=0{,}9$: $k\ge\frac{-4{,}605}{-0{,}1054}\approx43{,}7$, luego $k=44$. $\beta=0{,}999$: $k\ge4602{,}9$, luego $k=4603$.
\end{enumerate}
\resultado $G=10$ y $G\approx5{,}26$; $44$ y $4603$ pasos.
\comprobacion el «horizonte efectivo» de una media exponencial es del orden de $\frac{1}{1-\beta}$ pasos (10 y 1000); para reducir el sesgo al 1~\% hacen falta unas $\ln100\approx4{,}6$ veces más (la aproximación es mejor cuanto más cerca de~1 está $\beta$). Por eso Adam corrige el sesgo dividiendo por $1-\beta^{k}$.
\end{solucion}

\begin{solucion}{eje:5-armonica-dem}
\planteamiento acotar los bloques $\sum_{k=n+1}^{2n}\frac1k$ por debajo.
\pasos
\begin{enumerate}[label=\arabic*.]
  \item (a) $H_{2n}-H_n=\sum_{k=n+1}^{2n}\frac1k$ tiene $n$ sumandos, todos $\ge\frac{1}{2n}$; luego $H_{2n}-H_n\ge n\cdot\frac{1}{2n}=\frac12$.
  \item (b) Escribiendo $H_{2^{m}}=H_1+\sum_{j=0}^{m-1}\bigl(H_{2^{j+1}}-H_{2^{j}}\bigr)$ y aplicando (a) con $n=2^{j}$: $H_{2^{m}}\ge1+\frac m2$, que no está acotada.
  \item (c) Si $S_n\to S$, entonces $S_{2n}\to S$, y $\sum_{k=n+1}^{2n}a_k=S_{2n}-S_n\to0$.
  \item Para la serie armónica, $S_{2n}-S_n=H_{2n}-H_n\ge\frac12$ no tiende a $0$; por (c), la serie no converge.
\end{enumerate}
\resultado la serie armónica diverge, con tres argumentos.
\comprobacion (c) es un caso particular del criterio de Cauchy: en una serie convergente, los bloques de términos consecutivos lejanos suman cada vez menos.
\end{solucion}

\begin{solucion}{eje:5-cociente-dem}
\datos $a_n>0$, $\frac{a_{n+1}}{a_n}\to L$.
\planteamiento comparar con una serie geométrica de razón $r\in(L,1)$.
\pasos
\begin{enumerate}[label=\arabic*.]
  \item Caso $L<1$. Sea $r$ con $L<r<1$. Por la definición de límite con $\varepsilon=r-L$, existe $N$ con $\frac{a_{n+1}}{a_n}<r$ para todo $n\ge N$.
  \item Por inducción sobre $m$, $a_{N+m}\le a_Nr^{m}$: para $m=0$ es una igualdad, y si vale para $m$, $a_{N+m+1}<r\,a_{N+m}\le a_Nr^{m+1}$.
  \item La serie $\sum_{m\ge0}a_Nr^{m}$ converge (\cref{teo:geometrica}), luego $\sum_{n\ge N}a_n$ converge por comparación, y también $\sum a_n$, que solo añade un número finito de términos.
  \item Caso $L>1$ (o $L=\infty$). Existe $N$ con $\frac{a_{n+1}}{a_n}>1$ para $n\ge N$, luego $a_n\ge a_N>0$ para $n\ge N$, el término general no tiende a $0$ y la serie diverge.
\end{enumerate}
\resultado el \cref{teo:cociente} queda demostrado.
\comprobacion la hipótesis $r<1$ se usa en el paso~3; si $L=1$ no existe tal $r$, y por eso el criterio no decide.
\end{solucion}

\begin{solucion}{eje:5-resto}
\datos $f$ positiva y decreciente, $\int_1^{\infty}f$ convergente; $R_n=\sum_{k>n}f(k)$.
\planteamiento sumar las desigualdades $f(k+1)\le\int_k^{k+1}f\le f(k)$ de la demostración del \cref{teo:criterio-integral}.
\pasos
\begin{enumerate}[label=\arabic*.]
  \item Para $N>n$, sumando $\int_k^{k+1}f\le f(k)$ para $k=n+1,\dots,N$: $\int_{n+1}^{N+1}f\le\sum_{k=n+1}^{N}f(k)$.
  \item Sumando $f(k)\le\int_{k-1}^{k}f$ para $k=n+1,\dots,N$: $\sum_{k=n+1}^{N}f(k)\le\int_n^{N}f$.
  \item Haciendo $N\to\infty$ (todo converge): $\int_{n+1}^{\infty}f\le R_n\le\int_n^{\infty}f$.
  \item (b) Con $f(x)=x^{-2}$, $\int_m^{\infty}x^{-2}\dif x=\frac1m$, luego $\frac{1}{n+1}\le R_n\le\frac1n$.
  \item Con $n=10$: $S_{10}=1{,}54977$, y la suma está en $[S_{10}+\frac{1}{11},\,S_{10}+\frac{1}{10}]=[1{,}64068;\,1{,}64977]$. El punto medio, $1{,}64522$, la aproxima con error menor que $\frac12\bigl(\frac1{10}-\frac1{11}\bigr)\approx0{,}0045$.
\end{enumerate}
\resultado $\int_{n+1}^{\infty}f\le R_n\le\int_n^{\infty}f$; $\sum\frac{1}{k^{2}}\approx1{,}645$ con error menor que $0{,}005$.
\comprobacion el valor exacto, $\frac{\pi^{2}}{6}=1{,}64493$, está en el intervalo, y difiere del punto medio en $0{,}0003$.
\end{solucion}

\begin{solucion}{eje:5-euler}
\datos $c_n=H_n-\ln n$.
\planteamiento probar que $(c_n)$ es decreciente y acotada inferiormente, y aplicar el \cref{teo:monotona-acotada}.
\pasos
\begin{enumerate}[label=\arabic*.]
  \item $c_n-c_{n+1}=\ln(n+1)-\ln n-\frac{1}{n+1}=\int_n^{n+1}\frac{\dif x}{x}-\frac{1}{n+1}$.
  \item En $[n,n+1]$, $\frac1x\ge\frac{1}{n+1}$, luego $\int_n^{n+1}\frac{\dif x}{x}\ge\frac{1}{n+1}$ y $c_n-c_{n+1}\ge0$: la sucesión es decreciente.
  \item Por el \cref{cor:series-p}, $H_n\ge\ln(n+1)>\ln n$, luego $c_n>0$.
  \item Decreciente y acotada inferiormente por $0$: converge (\cref{teo:monotona-acotada}).
\end{enumerate}
\resultado $(c_n)$ converge; su límite $\gamma$ cumple $0\le\gamma<c_1=1$.
\comprobacion $c_{10}=2{,}92897-2{,}30259=0{,}62638$, por encima de $\gamma\approx0{,}57722$, como corresponde a una sucesión decreciente.
\end{solucion}

\begin{solucion}{eje:5-heron}
\datos $a_1=2$, $a_{n+1}=\frac12\bigl(a_n+\frac{2}{a_n}\bigr)$.
\planteamiento cota inferior por la desigualdad de las medias, monotonía por diferencia, y paso al límite en la recurrencia.
\pasos
\begin{enumerate}[label=\arabic*.]
  \item Por inducción, $a_n>0$ para todo $n$. Para $n\ge1$, por la desigualdad entre las medias aritmética y geométrica de $a_n$ y $\frac{2}{a_n}$: $a_{n+1}\ge\sqrt{a_n\cdot\frac{2}{a_n}}=\sqrt2$. Además $a_1=2\ge\sqrt2$.
  \item $a_n-a_{n+1}=\frac{a_n}{2}-\frac{1}{a_n}=\frac{a_n^{2}-2}{2a_n}\ge0$, por el paso~1. La sucesión es decreciente.
  \item Decreciente y acotada inferiormente por $\sqrt2$: converge a un $\ell\ge\sqrt2>0$.
  \item Pasando al límite en la recurrencia (álgebra de límites, $\ell\neq0$): $\ell=\frac12\bigl(\ell+\frac2\ell\bigr)$, de donde $\ell^{2}=2$ y $\ell=\sqrt2$.
  \item $a_2=\frac32=1{,}5$; $a_3=\frac12\bigl(\frac32+\frac43\bigr)=\frac{17}{12}\approx1{,}416667$; $a_4=\frac12\bigl(\frac{17}{12}+\frac{24}{17}\bigr)=\frac{577}{408}\approx1{,}414216$.
\end{enumerate}
\resultado $a_n\to\sqrt2$; $a_4$ difiere de $\sqrt2=1{,}414214$ en unas $2\cdot10^{-6}$.
\comprobacion es el método de Newton aplicado a $x^{2}-2=0$; el número de cifras correctas se duplica aproximadamente en cada paso: errores $0{,}086$, $0{,}0025$ y $2{,}1\cdot10^{-6}$.
\end{solucion}

% ---------------------------------------------------------------------
\section{Tema 6}

\begin{solucion}{eje:6-vf}
\planteamiento \cref{lem:abel} y \cref{teo:radio}.
\pasos
\begin{enumerate}[label=(\alph*)]
  \item \textbf{Verdadera}: $\abs{-2}<\abs3$, y el \cref{lem:abel} da convergencia absoluta en $x=-2$.
  \item \textbf{Falsa}: $\sum\frac{(-1)^{n}x^{n}}{n\,3^{n}}$ converge en $x=3$ (serie armónica alternada) y diverge en $x=-3$ (serie armónica).
  \item \textbf{Verdadera}: en $x=c$ todos los términos con $n\ge1$ son nulos.
  \item \textbf{Falsa} tal como está: $\ln5=\ln(1+4)$ y $x=4$ está fuera del intervalo $(-1,1]$. Sí puede usarse con un truco: $\ln5=-\ln\frac15=-\ln\bigl(1+(-\frac45)\bigr)$, con $x=-0{,}8$ dentro del intervalo.
\end{enumerate}
\resultado (a) V; (b) F; (c) V; (d) F.
\comprobacion (b) muestra que en los extremos del intervalo de convergencia puede ocurrir cualquier cosa.
\end{solucion}

\begin{solucion}{eje:6-radio}
\planteamiento $R=\lim\abs{a_n/a_{n+1}}$ (\cref{teo:radio}) y estudio de los extremos con el \cref{tema:5}.
\pasos
\begin{enumerate}[label=(\alph*)]
  \item $\frac{(n+1)^{2}}{n^{2}}\to1$: $R=1$. En $x=\pm1$, $\sum\frac{(\pm1)^{n}}{n^{2}}$ converge absolutamente. Dominio $[-1,1]$.
  \item $\frac{n}{n+1}\to1$: $R=1$. En $x=\pm1$, $n(\pm1)^{n}\not\to0$. Dominio $(-1,1)$.
  \item $\frac{3^{n+1}(n+1)}{3^{n}n}\to3$: $R=3$, centro $2$. En $x=5$: $\sum\frac1n$, diverge; en $x=-1$: $\sum\frac{(-1)^{n}}{n}$, converge. Dominio $[-1,5)$.
  \item Es la geométrica $\sum\bigl(-\frac{x^{2}}{4}\bigr)^{n}$, que converge si y solo si $\frac{x^{2}}{4}<1$: $R=2$. En $x=\pm2$, $\sum(-1)^{n}$ diverge. Dominio $(-2,2)$, con suma $\frac{4}{4+x^{2}}$.
\end{enumerate}
\resultado (a) $[-1,1]$; (b) $(-1,1)$; (c) $[-1,5)$; (d) $(-2,2)$.
\comprobacion en (d) la fórmula del \cref{teo:radio} no se aplica directamente, porque los coeficientes impares son nulos; tratarla como serie en $x^{2}$ evita el problema.
\end{solucion}

\begin{solucion}{eje:6-e2x}
\planteamiento sustituir $2x$ en las series del \cref{teo:exp} y del seno.
\pasos
\begin{enumerate}[label=\arabic*.]
  \item (a) $e^{2x}=\sum_{n\ge0}\frac{(2x)^{n}}{n!}=\sum_{n\ge0}\frac{2^{n}}{n!}x^{n}$. Como $f^{(n)}(x)=2^{n}e^{2x}$, $\frac{f^{(n)}(0)}{n!}=\frac{2^{n}}{n!}$, de acuerdo con el \cref{cor:coeficientes}. Con $x=1$: $e^{2}=\sum_{n\ge0}\frac{2^{n}}{n!}$.
  \item (b) $\sin2x=\sum_{n\ge0}\frac{(-1)^{n}(2x)^{2n+1}}{(2n+1)!}=2x-\frac43x^{3}+\frac{4}{15}x^{5}-\cdots$.
  \item El coeficiente de $x^{3}$ es $\frac{f'''(0)}{3!}=-\frac43$, luego $f'''(0)=-8$; directamente, $(\sin2x)'''=-8\cos2x$, que vale $-8$ en $0$.
  \item Con $x=1$: $\sin2=\sum_{n\ge0}\frac{(-1)^{n}2^{2n+1}}{(2n+1)!}$.
\end{enumerate}
\resultado $e^{2}=\sum\frac{2^{n}}{n!}$; $f'''(0)=-8$; $\sin2=\sum\frac{(-1)^{n}2^{2n+1}}{(2n+1)!}$.
\comprobacion con cinco términos, $1+2+2+\frac43+\frac23=7$, y $e^{2}\approx7{,}389$; la serie converge rápidamente por el factorial.
\end{solucion}

\begin{solucion}{eje:6-taylor-basicas}
\planteamiento calcular $f^{(n)}(0)$ y aplicar la \cref{def:serie-taylor}.
\pasos
\begin{enumerate}[label=(\alph*)]
  \item $f(x)=(1+x)^{-1}$: $f^{(n)}(x)=(-1)^{n}n!\,(1+x)^{-n-1}$, $\frac{f^{(n)}(0)}{n!}=(-1)^{n}$. Serie: $\sum_{n\ge0}(-1)^{n}x^{n}=1-x+x^{2}-\cdots$.
  \item $\cosh'=\sinh$, $\sinh'=\cosh$, con $\cosh0=1$ y $\sinh0=0$: las derivadas pares valen~1 y las impares~0 en $0$. Serie: $\sum_{n\ge0}\frac{x^{2n}}{(2n)!}=1+\frac{x^{2}}{2}+\frac{x^{4}}{24}+\cdots$.
  \item $f(x)=\ln(1+x)$: $f'=(1+x)^{-1}$, y por (a), $f^{(n)}(x)=(-1)^{n-1}(n-1)!\,(1+x)^{-n}$ para $n\ge1$; $\frac{f^{(n)}(0)}{n!}=\frac{(-1)^{n-1}}{n}$. Serie: $x-\frac{x^{2}}{2}+\frac{x^{3}}{3}-\cdots$.
\end{enumerate}
\resultado las tres series anteriores.
\comprobacion (b): $\frac12\bigl(\sum\frac{x^{n}}{n!}+\sum\frac{(-x)^{n}}{n!}\bigr)$ cancela los términos impares y conserva los pares. (c) coincide con el \cref{teo:ln}.
\end{solucion}

\begin{solucion}{eje:6-derivar-geom}
\planteamiento derivar término a término la serie geométrica (\cref{teo:termino-a-termino}) e integrar la serie de $\frac1x$.
\pasos
\begin{enumerate}[label=\arabic*.]
  \item (a) $\sum_{n\ge1}nx^{n-1}=\frac{1}{(1-x)^{2}}$ y $\sum_{n\ge2}n(n-1)x^{n-2}=\frac{2}{(1-x)^{3}}$ en $(-1,1)$. Con $x=\frac12$: $\frac{1}{1/4}=4$ y $\frac{2}{1/8}=16$.
  \item (b) $\sum_{k\ge1}k\,2^{-k}=\frac12\sum_{k\ge1}k\bigl(\frac12\bigr)^{k-1}=\frac12\cdot4=2$. Con $x=1-p\in(0,1)$: $\sum_{k\ge1}k(1-p)^{k-1}p=\frac{p}{p^{2}}=\frac1p$.
  \item (c) Para $\abs{1-x}<1$, es decir, $0<x<2$: $\frac1x=\frac{1}{1-(1-x)}=\sum_{n\ge0}(1-x)^{n}$. Integrando de $1$ a $x$: $\ln x=\sum_{n\ge0}\eval{-\frac{(1-t)^{n+1}}{n+1}}{1}{x}=-\sum_{n\ge1}\frac{(1-x)^{n}}{n}=\sum_{n\ge1}\frac{(-1)^{n+1}}{n}(x-1)^{n}$.
\end{enumerate}
\resultado $4$, $16$; $\sum k\,2^{-k}=2$; $\sum k(1-p)^{k-1}p=\frac1p$; $\ln x=\sum\frac{(-1)^{n+1}}{n}(x-1)^{n}$ en $(0,2)$.
\comprobacion $\frac1p$ es la media de la distribución geométrica, el número esperado de intentos hasta el primer éxito. (c) es el \cref{teo:ln} con $x$ sustituido por $x-1$.
\end{solucion}

\begin{solucion}{eje:6-productos}
\planteamiento dividir, multiplicar o componer series conocidas, conservando los términos hasta el grado necesario.
\pasos
\begin{enumerate}[label=(\alph*)]
  \item $\frac{\sin x}{x}=\frac1x\bigl(x-\frac{x^{3}}{6}+\frac{x^{5}}{120}-\cdots\bigr)=1-\frac{x^{2}}{6}+\frac{x^{4}}{120}-\cdots$.
  \item $\bigl(1+x+\frac{x^{2}}{2}+\frac{x^{3}}{6}+\frac{x^{4}}{24}\bigr)\bigl(x-\frac{x^{3}}{6}+\frac{x^{5}}{120}\bigr)$. Por grados: $x$; $x^{2}$; $x^{3}\bigl(\frac12-\frac16\bigr)=\frac{x^{3}}{3}$; $x^{4}\bigl(\frac16-\frac16\bigr)=0$; $x^{5}\bigl(\frac{1}{24}-\frac{1}{12}+\frac{1}{120}\bigr)=-\frac{x^{5}}{30}$. Resultado: $x+x^{2}+\frac{x^{3}}{3}-\frac{x^{5}}{30}+\cdots$.
  \item Sustituyendo $-x^{2}$ en la exponencial: $e^{-x^{2}}=1-x^{2}+\frac{x^{4}}{2}-\frac{x^{6}}{6}+\cdots$. Integrando término a término: $\int_0^{x}e^{-t^{2}}\dif t=x-\frac{x^{3}}{3}+\frac{x^{5}}{10}-\frac{x^{7}}{42}+\cdots$.
\end{enumerate}
\resultado las series anteriores.
\comprobacion en (b), $e^{x}\sin x=\operatorname{Im}e^{(1+i)x}$, y el coeficiente de $x^{4}$ es $\frac{\operatorname{Im}(1+i)^{4}}{4!}=\frac{\operatorname{Im}(-4)}{24}=0$. En (c), con $x=1$ los cuatro términos dan $0{,}7429$, frente a $\int_0^{1}e^{-t^{2}}\dif t\approx0{,}7468$.
\end{solucion}

\begin{solucion}{eje:6-limites}
\planteamiento sustituir cada función por los primeros términos de su serie y simplificar.
\pasos
\begin{enumerate}[label=(\alph*)]
  \item $e^{x}-1-x=\frac{x^{2}}{2}+\frac{x^{3}}{6}+\cdots$; dividido por $x^{2}$: $\frac12+\frac x6+\cdots\to\frac12$.
  \item $x-\sin x=\frac{x^{3}}{6}-\frac{x^{5}}{120}+\cdots$; dividido por $x^{3}$: $\to\frac16$.
  \item $\ln(1+x)-x=-\frac{x^{2}}{2}+\frac{x^{3}}{3}-\cdots$; dividido por $x^{2}$: $\to-\frac12$.
  \item $1-\cos x=\frac{x^{2}}{2}-\frac{x^{4}}{24}+\cdots$; dividido por $x^{2}$: $\to\frac12$.
\end{enumerate}
\resultado (a) $\frac12$; (b) $\frac16$; (c) $-\frac12$; (d) $\frac12$.
\comprobacion (b) por L'Hôpital tres veces: $\frac{1-\cos x}{3x^{2}}\to\frac{\sin x}{6x}\to\frac16$.
\end{solucion}

\begin{solucion}{eje:6-cos1}
\datos $\cos1=\sum_{n\ge0}\frac{(-1)^{n}}{(2n)!}$; cinco decimales exigen error menor que $5\cdot10^{-6}$.
\planteamiento la serie es alternada con términos decrecientes; el error es menor que el primer término omitido (\cref{teo:leibniz}).
\pasos
\begin{enumerate}[label=\arabic*.]
  \item Términos: $1$, $\frac12$, $\frac{1}{24}=0{,}0416667$, $\frac{1}{720}=0{,}0013889$, $\frac{1}{40\,320}=2{,}48\cdot10^{-5}$, $\frac{1}{3\,628\,800}=2{,}76\cdot10^{-7}$.
  \item Si se para en $\frac{1}{720}$, el error puede llegar a $2{,}48\cdot10^{-5}>5\cdot10^{-6}$: no basta.
  \item Con los cinco primeros términos (hasta $\frac{1}{8!}$): $1-0{,}5+0{,}0416667-0{,}0013889+0{,}0000248=0{,}5403026$, con error menor que $2{,}76\cdot10^{-7}$.
\end{enumerate}
\resultado $\cos1\approx0{,}54030$; bastan cinco términos.
\comprobacion $\cos1=0{,}5403023$; la suma parcial difiere en $2{,}7\cdot10^{-7}$, por debajo de la cota.
\end{solucion}

\begin{solucion}{eje:6-aproximar}
\planteamiento \cref{teo:binomial}, \cref{teo:ln} con la cota de Leibniz y \cref{teo:e-irracional}.
\pasos
\begin{enumerate}[label=\arabic*.]
  \item (a) $\binom{1/2}{2}=\frac{(1/2)(-1/2)}{2}=-\frac18$, $\binom{1/2}{3}=\frac{(1/2)(-1/2)(-3/2)}{6}=\frac{1}{16}$: $(1+x)^{1/2}=1+\frac x2-\frac{x^{2}}{8}+\frac{x^{3}}{16}-\cdots$.
  \item $\binom{-1/2}{2}=\frac{(-1/2)(-3/2)}{2}=\frac38$, $\binom{-1/2}{3}=\frac{(-1/2)(-3/2)(-5/2)}{6}=-\frac{5}{16}$: $(1+x)^{-1/2}=1-\frac x2+\frac{3x^{2}}{8}-\frac{5x^{3}}{16}+\cdots$.
  \item Con $x=0{,}1$: $\sqrt{1{,}1}\approx1+0{,}05-0{,}00125+0{,}0000625=1{,}0488125$.
  \item (b) $\ln1{,}2\approx0{,}2-\frac{0{,}04}{2}+\frac{0{,}008}{3}-\frac{0{,}0016}{4}=0{,}1822667$. Es alternada y decreciente: error $\le\frac{0{,}2^{5}}{5}=6{,}4\cdot10^{-5}$.
  \item (c) Se busca $q$ con $\frac{1}{q!\,q}<10^{-6}$: $q=8$ da $3{,}1\cdot10^{-6}$; $q=9$ da $3{,}06\cdot10^{-7}$. Basta $s_9$, es decir, los diez términos $\frac{1}{0!},\dots,\frac{1}{9!}$.
\end{enumerate}
\resultado $\sqrt{1{,}1}\approx1{,}048813$; $\ln1{,}2\approx0{,}18227$ con error $\le6{,}4\cdot10^{-5}$; $s_9$.
\comprobacion $\sqrt{1{,}1}=1{,}0488088$ (error $3{,}7\cdot10^{-6}$, del orden del término siguiente, $\frac{5}{128}\cdot10^{-4}$); $\ln1{,}2=0{,}1823216$ (error $5{,}5\cdot10^{-5}$, dentro de la cota).
\end{solucion}

\begin{solucion}{eje:6-seno-dem}
\planteamiento como en el \cref{teo:exp}: calcular las derivadas en $0$ y acotar el resto de Lagrange.
\pasos
\begin{enumerate}[label=\arabic*.]
  \item Las derivadas de $\sin$ son, cíclicamente, $\cos,-\sin,-\cos,\sin$; en $0$ valen $0,1,0,-1,0,1,\dots$ para $\sin^{(0)},\sin^{(1)},\dots$ Luego los coeficientes pares son nulos y el de $x^{2n+1}$ es $\frac{(-1)^{n}}{(2n+1)!}$: la serie del enunciado es la de Taylor.
  \item Todas las derivadas de $\sin$ están acotadas por~1 en valor absoluto. Por el \cref{teo:taylor-lagrange}, $\abs{R_n(x)}\le\frac{\abs x^{n+1}}{(n+1)!}$.
  \item $\frac{\abs x^{m}}{m!}\to0$ por ser el término general de una serie convergente (\cref{ej:cociente}); luego $R_n(x)\to0$ y la serie converge a $\sin x$ (\cref{prop:taylor-converge}).
  \item El mismo argumento, con las derivadas de $\cos$ en $0$ iguales a $1,0,-1,0,\dots$, prueba la serie del coseno.
  \item Ambas series tienen radio $\infty$; por el \cref{teo:termino-a-termino}, $\frac{d}{dx}\sum\frac{(-1)^{n}x^{2n+1}}{(2n+1)!}=\sum\frac{(-1)^{n}(2n+1)x^{2n}}{(2n+1)!}=\sum\frac{(-1)^{n}x^{2n}}{(2n)!}$.
\end{enumerate}
\resultado las series convergen a $\sin x$ y $\cos x$ en todo $\R$, y $(\sin)'=\cos$ se lee en las series.
\comprobacion derivando la serie del coseno se obtiene $-\sum_{n\ge1}\frac{(-1)^{n-1}x^{2n-1}}{(2n-1)!}=-\sin x$.
\end{solucion}

\begin{solucion}{eje:6-arctan}
\planteamiento integrar una identidad exacta con resto, como en el \cref{teo:ln}.
\pasos
\begin{enumerate}[label=\arabic*.]
  \item (a) La identidad es la suma geométrica finita con razón $-t^{2}$. Integrando entre $0$ y $x$: $\arctan x=\sum_{k=0}^{n-1}\frac{(-1)^{k}x^{2k+1}}{2k+1}+E_n(x)$, con $E_n(x)=(-1)^{n}\int_0^{x}\frac{t^{2n}}{1+t^{2}}\dif t$.
  \item Como $1+t^{2}\ge1$: $\abs{E_n(x)}\le\Bigl|\int_0^{x}t^{2n}\dif t\Bigr|=\frac{\abs x^{2n+1}}{2n+1}\le\frac{1}{2n+1}$ si $\abs x\le1$. Luego $E_n(x)\to0$ en $[-1,1]$.
  \item (b) Con $x=1$: $\frac\pi4=\arctan1=\sum_{k\ge0}\frac{(-1)^{k}}{2k+1}=1-\frac13+\frac15-\cdots$.
  \item (c) Con $n$ términos el error es menor o igual que $\frac{1}{2n+1}$; para que sea menor que $10^{-3}$ hace falta $2n+1>1000$, es decir, $n\ge500$.
\end{enumerate}
\resultado la serie converge a $\arctan x$ en $[-1,1]$; $\frac\pi4=1-\frac13+\frac15-\cdots$; 500 términos.
\comprobacion con 10 términos la suma es $0{,}7605$ y $\frac\pi4=0{,}7854$: error $0{,}025<\frac{1}{21}$. La serie de Leibniz converge muy despacio; en la práctica se usan fórmulas como $\frac\pi4=4\arctan\frac15-\arctan\frac{1}{239}$.
\end{solucion}

\begin{solucion}{eje:6-binomial-dem}
\datos $f(x)=\sum_{n\ge0}\binom anx^{n}$, $\abs x<1$ (radio $R\ge1$).
\planteamiento probar la ecuación diferencial $(1+x)f'=af$ con la identidad de coeficientes y deducir $f=(1+x)^{a}$.
\pasos
\begin{enumerate}[label=\arabic*.]
  \item Identidad: $\binom{a}{n+1}=\binom an\frac{a-n}{n+1}$, luego $(n+1)\binom{a}{n+1}=(a-n)\binom an$ y $(n+1)\binom{a}{n+1}+n\binom an=a\binom an$.
  \item Por el \cref{teo:termino-a-termino}, $f'(x)=\sum_{n\ge1}n\binom anx^{n-1}=\sum_{n\ge0}(n+1)\binom{a}{n+1}x^{n}$, y $xf'(x)=\sum_{n\ge0}n\binom anx^{n}$.
  \item Sumando: $(1+x)f'(x)=\sum_{n\ge0}\Bigl[(n+1)\binom{a}{n+1}+n\binom an\Bigr]x^{n}=a\sum_{n\ge0}\binom anx^{n}=af(x)$.
  \item (b) Sea $g(x)=f(x)(1+x)^{-a}$ en $(-1,1)$. $g'(x)=(1+x)^{-a-1}\bigl[(1+x)f'(x)-af(x)\bigr]=0$.
  \item $g$ es constante en el intervalo $(-1,1)$ e igual a $g(0)=f(0)=\binom a0=1$. Luego $f(x)=(1+x)^{a}$.
\end{enumerate}
\resultado $(1+x)^{a}=\sum_{n\ge0}\binom anx^{n}$ para $\abs x<1$.
\comprobacion con $a=-1$, $\binom{-1}{n}=(-1)^{n}$ y se recupera la serie geométrica $\frac{1}{1+x}=\sum(-1)^{n}x^{n}$; con $a=2$, la serie se corta: $1+2x+x^{2}$.
\end{solucion}

\begin{solucion}{eje:6-limite-exp}
\datos $x\in\R$ fijo.
\planteamiento escribir $\bigl(1+\frac xn\bigr)^{n}=\exp\bigl(n\ln(1+\frac xn)\bigr)$ y usar la derivada del logaritmo en~$1$.
\pasos
\begin{enumerate}[label=\arabic*.]
  \item Si $x=0$, la sucesión es constante e igual a $1=e^{0}$.
  \item Si $x\neq0$, para $n>\abs x$ la base es positiva y $n\ln\bigl(1+\frac xn\bigr)=x\cdot\frac{\ln(1+h)}{h}$ con $h=\frac xn\to0$.
  \item Como $\frac{\ln(1+h)}{h}\to(\ln)'(1)=1$, el exponente tiende a $x$, y por la continuidad de la exponencial, $\bigl(1+\frac xn\bigr)^{n}\to e^{x}$.
  \item Con $x=-1$: $\bigl(1-\frac1n\bigr)^{n}\to e^{-1}\approx0{,}3679$.
  \item En $n$ extracciones independientes con reemplazamiento entre $n$ elementos, la probabilidad de que un elemento fijo no salga en una extracción es $1-\frac1n$, y en ninguna, $\bigl(1-\frac1n\bigr)^{n}\to\frac1e$.
\end{enumerate}
\resultado $\lim\bigl(1+\frac xn\bigr)^{n}=e^{x}$; $\lim\bigl(1-\frac1n\bigr)^{n}=\frac1e$.
\comprobacion con $n=100$, $(0{,}99)^{100}\approx0{,}3660$. Es el porcentaje de datos que quedan fuera de una muestra \emph{bootstrap}: en torno al 36,8~\%.
\end{solucion}

\begin{solucion}{eje:6-no-analitica}
\datos $f(x)=e^{-1/x^{2}}$ para $x\neq0$, $f(0)=0$.
\planteamiento (a) inducción sobre $n$ con la regla de la cadena; (b) inducción usando la definición de derivada en $0$; (c) comparar la serie de Taylor con $f$.
\pasos
\begin{enumerate}[label=\arabic*.]
  \item (a) Para $n=0$, $P_0=1$. Si $f^{(n)}(x)=P_n(1/x)e^{-1/x^{2}}$ para $x\neq0$, derivando con $y=1/x$, $\frac{dy}{dx}=-y^{2}$ y $\frac{d}{dx}e^{-1/x^{2}}=2y^{3}e^{-1/x^{2}}$: $f^{(n+1)}(x)=\bigl[-y^{2}P_n'(y)+2y^{3}P_n(y)\bigr]e^{-1/x^{2}}$, y $P_{n+1}(y)=2y^{3}P_n(y)-y^{2}P_n'(y)$ es un polinomio.
  \item (b) Para todo $k\in\N_0$, $y^{k}e^{-y^{2}}\to0$ cuando $\abs y\to\infty$ (la exponencial domina a las potencias); por linealidad, $yP(y)e^{-y^{2}}\to0$ para todo polinomio $P$.
  \item Inducción: $f^{(0)}(0)=0$. Si $f^{(n)}(0)=0$, entonces $f^{(n+1)}(0)=\lim_{h\to0}\frac{f^{(n)}(h)}{h}=\lim_{h\to0}\frac1hP_n\bigl(\tfrac1h\bigr)e^{-1/h^{2}}=\lim_{\abs y\to\infty}y\,P_n(y)e^{-y^{2}}=0$.
  \item (c) Todos los coeficientes de Taylor en $0$ son nulos: la serie es la serie nula, que converge en todo $\R$ con suma $0$. Pero $f(x)>0$ para $x\neq0$, así que la serie solo coincide con $f$ en $x=0$.
\end{enumerate}
\resultado $f$ es infinitamente derivable, todas sus derivadas en $0$ son nulas, y su serie de Taylor no representa a $f$ en ningún entorno de~$0$.
\comprobacion por la \cref{prop:taylor-converge}, el resto $R_nf(x)=f(x)$ no tiende a $0$ para $x\neq0$. Las funciones que sí coinciden con su serie de Taylor en un entorno de cada punto se llaman \emph{analíticas}; $e^{x}$, $\sin$, $\cos$ y $\ln(1+x)$ lo son.
\end{solucion}
